% Esta es la Plantilla UNAL en LaTeX
\documentclass[10pt,spanish,fleqn,openany,twoside,letterpaper]{book}

%Muestra los márgenes del documento para evitar Warnings
%Para activar la siguiente línea quite el simbolo % 
%\usepackage[showframe]{geometry}

%Formato de fuentes bibliográficas
%Use el estilo bibliográfico que sea pertinente según el área de estudio APA, IEEE, etc

%Usando el paquete BibLaTeX
%Cita normal con \cite[página]{} y cita con paréntesis \parencite[página]{}
% Configuración de BibLaTeX
\usepackage[spanish]{babel}
\usepackage[T1]{fontenc}
\usepackage[utf8]{inputenc}
\usepackage{chngcntr}
\usepackage[backend=biber,style=apa,maxcitenames=2,maxbibnames=99,giveninits=true,uniquename=false]{biblatex}
\addbibresource{archivos_bib/Referencias.bib}
\addbibresource{archivos_bib/Bibliografia_unificada.bib}
% Cambiar el idioma de las referencias bibliográficas a español
\DefineBibliographyExtras{spanish}{\def\keywordname{Referencias}}
%\DefineBibliographyStrings{spanish}{%
%  andothers = {et\addabbrvspace al\adddot},
%  andmore = {et\addabbrvspace al\adddot},
%}

% Personalizar el formato de las citas y la bibliografía
%\DeclareNameAlias{sortname}{family-given}
%\DeclareDelimFormat{multinamedelim}{\addcomma\space}
%\DeclareDelimFormat{finalnamedelim}{\addcomma\space\&\space}
%\DeclareFieldFormat{titlecase}{\MakeSentenceCase*{#1}}
%\DeclareFieldFormat[article,inbook,incollection,inproceedings,patent,thesis,unpublished]{title}{\titlecase{#1}}
%\DeclareFieldFormat{journaltitlecase}{\titlecase{#1}}
%\DeclareFieldFormat{pages}{#1}
%\DeclareFieldFormat{volume}{\mkbibbold{#1}}
%\renewbibmacro{in:}{}
%\AtEveryBibitem{\clearfield{month}}

%Usando el paquete Natbib
%Cita normal \cite[página]{} y cita con paréntesis \citep[página]{}
%\bibpunct{[}{]}{;}{\&}{.}{}
%\bibliographystyle{dtvstyle}

%Idioma del documento
%\usepackage{natbib}
%Use main para el idioma principal del documento
\usepackage[spanish]{babel}

% Carácteres especiales
\usepackage{fontenc}
\usepackage{comment}

% Evita ligadura li & fl
\usepackage{microtype}
\DisableLigatures{encoding = *, family = *}

% Otros paquetes de tablas y colores avanzados
\usepackage{amsmath,graphicx,rotating,float,multirow}
\usepackage{longtable}
\usepackage{pgfgantt}
\setlength{\LTcapwidth}{6in}
\usepackage{epsfig,epic,eepic,threeparttable,amscd,here,lscape,tabularx,subfigure}
\usepackage{tabu,array}
\usepackage[rgb]{xcolor}
\let\cleardoublepage\clearpage
% Permite ver y configurar los parámetros de la página
\usepackage{layout}
%Hyperref permite ver las secciones del texto
\usepackage[hidelinks]{hyperref}
\usepackage{enumitem} % Ensure this package is included for custom enumerate labels
\usepackage[version=4]{mhchem}
\usepackage{siunitx}
\usepackage{chemformula}
\usepackage{chemfig}
\usepackage{chemmacros}
\usepackage{csquotes}

%Genera los comandos de la página de autoría
\newcommand{\studentname}{}
\newcommand{\submissiondate}{}
\newcommand{\academictitle}{}
\newcommand{\resgroupone}{}
\newcommand{\resgrouptwo}{}
\newcommand{\researchtopic}{}
\newcommand{\thesisname}{}
\newcommand{\director}{}
\newcommand{\codirector}{}
\newcommand{\codirectora}{}
\newcommand{\issuedate}{}
\newcommand{\palabrasclave}{}
\newcommand{\keywords}{}
\newcommand{\schlusselworter}{}
\newcommand{\palavraschave}{}
\newcommand{\sede}{}
\newcommand{\department}{}
\newcommand{\faculty}{}

%Información de la tesis
%Diligenciar aquí los datos para su carga automática donde se requiera en el documento
\renewcommand{\studentname}{John Jairo Leal Gómez}
\renewcommand{\thesisname}{Sistema Híbrido para Predicción de Emisiones de CO$_2$ y CH$_4$ en la Bioconversión de Residuos con \textit{Hermetia illucens}}
\renewcommand{\issuedate}{2025}
\renewcommand{\submissiondate}{Abril}
\renewcommand{\director}{Director Dr. Luis Octavio González Salcedo}
\renewcommand{\codirector}{Codirector Dr. Luis Fernando Mejía Rodríguez }
\renewcommand{\codirectora}{Codirector Dra. Sorany Milena Barrientos Grajales }
\renewcommand{\academictitle}{Doctor (PhD)  }
\renewcommand{\resgroupone}{Grupo A (Sigla Grupo Investigación 01) }
\renewcommand{\resgrouptwo}{Grupo B (Sigla Grupo Investigación 02) }
\renewcommand{\researchtopic}{Línea}
\renewcommand{\sede}{Sede Palmira} 
\renewcommand{\department}{Doctorado en Estudios Ambientales}
\renewcommand{\faculty}{Facultad de Ingeniería y Administración}

%Palabras clave del documento - Tener presente los Theasurus https://www.thesaurus.com/
%Disponible en 3 idiomas aunque se puede extender a francés o otro idioma
\renewcommand{\palabrasclave}{Modelamiento matemático, \textit{Hermetia Illucens}, IOT, Algoritmos de Aprendizaje Automático} 
\renewcommand{\keywords}{Use keywords available in Theasaurus}
%\renewcommand{\schlusselworter}{}
%\renewcommand{\palavraschave}{}

% Estilo de los encabezados y pies de página
\usepackage{fancyhdr}
\fancyhf{}%
\pagestyle{fancyplain}
\textheight22.5cm \topmargin0cm \textwidth16.5cm \headheight22pt
\oddsidemargin0.5cm \evensidemargin-0.5cm%
\fancypagestyle{plain}{
\fancyhead[RO,LE]{}
\fancyhead[RE,LO]{\small \textbf{\thesisname}}
\fancyfoot[CO,CE]{\thepage}
}
\pagestyle{fancy}
\fancyhf{}%
\renewcommand{\chaptermark}[1]{\markboth{\thechapter.\; #1}{}}
\renewcommand{\sectionmark}[1]{\markright{\thesection.\; #1}{}}
\fancyhead[LO,RE]{\leftmark}
\fancyhead[RO,LE]{\rightmark}
\fancyfoot[CO,CE]{\thepage}
\thispagestyle{fancy}%

\usepackage{titlesec}
% Permite personalizar los títulos de sección y de capítulos
% hang lo deja en el mismo renglón, display lo despliega
% Elimina el "Capitulo" y deja solo el número
\titleformat{\chapter}[hang]
  {\sffamily\Huge\bfseries}{\thechapter}{0.5cm}{\sffamily\Huge}
\titleformat{\section}[hang]{\sffamily\LARGE}{\thesection}{0.5cm}{}
\titleformat{\subsection}[hang]{\sffamily\Large}{\thesubsection}{0.5cm}{}
\titleformat{\subsubsection}[hang]{\sffamily\large}{\thesubsubsection}{0.5cm}{}
\titleformat{\paragraph}[runin]{\sffamily\normalsize}{}{}{\emph}
% Reducción de espaciado entre títulos y contenido
\titlespacing{\chapter}{0pt}{-10pt}{10pt} % Reduce espacio antes y después de capítulo
\titlespacing{\section}{0pt}{8pt}{4pt}    % Reduce espacio antes y después de sección
\titlespacing{\subsection}{0pt}{6pt}{3pt} % Reduce espacio antes y después de subsección
\titlespacing{\subsubsection}{0pt}{4pt}{2pt} % Reduce espacio antes y después de subsubsección
\counterwithout{subsubsection}{subsection}  % Elimina la dependencia de subsection
\counterwithin{subsubsection}{chapter}      % Vincula subsubsection al chapter

\renewcommand{\thesubsubsection}{\thechapter.\arabic{subsubsection}}
%Coloca anexo o apéndice en la Tabla de contenido
\usepackage[toc,page]{appendix}


% Configuración de las páginas en twoside-mode
% Permite ver y configurar los parámetros de la página
\setlength{\voffset}{-0.25in}
\setlength{\headwidth}{467pt}
\setlength{\headheight}{22pt}
\setlength{\oddsidemargin}{0pt}
\setlength{\evensidemargin}{0pt}
\setlength{\marginparwidth}{0pt}
\setlength{\marginparsep}{0pt}
\setlength{\parskip}{2em}
\setlength{\footskip}{20pt}
\setlength{\textheight}{650pt}
\setlength{\textwidth}{467pt}
\setlength{\headsep}{5pt}
\setlength{\parindent}{0pt}
\setlength{\baselineskip}{10pt plus 5pt minus 5pt}
\renewcommand{\theequation}{\thechapter-\arabic{equation}}
\renewcommand{\thefigure}{\textbf{\thechapter-\arabic{figure}}}
\renewcommand{\thetable}{\textbf{\thechapter-\arabic{table}}}


%Define la distancia de la primera linea de un parrafo a la margen
\parindent0cm 

%Espacio entre lineas
\renewcommand{\baselinestretch}{1}

%Para rotar texto, objetos y tablas seite.
\usepackage{rotating}

%Permite incluir mecanismos y reacciones químicas
\usepackage{tikz}
\usetikzlibrary{positioning}
\usepackage{chemformula}
\usepackage{chemfig}

\usetikzlibrary{calc,arrows.meta}% per right to e left to
\tikzset{
myedge/.style={->, -{Latex[#1]}}
}

%Fuente de la presentación Ancizar Sans UNAL
%Para usar este compilado en Overleaf se debe usar el compilador XeLaTeX o LuaLaTeX!!
%Menu -> Compiler -> XeLaTeX o LuaLaTeX
%La siguiente línea debe comentarse si desea compilar con pdfLaTeX
%\RequireXeTeX

% Definición de la fuente Ancizar Sans
\newif\ifxetexorluatex

\ifxetexorluatex
  \usepackage{fontspec}
  \usefonttheme{serif}
  \setmainfont{AncizarSans}[Path=./AncizarSans/,Scale=1,Extension=.otf,UprightFont=*-Regular,BoldFont=*-Bold,ItalicFont=*-Italic,BoldItalicFont=*-BoldItalic]
\else
  % Si se compila con pdfLaTeX, cargar la fuente apropiada aquí
  \usepackage[T1]{fontenc}
\fi
% Metadatos del documento
\AtBeginDocument{%
	\hypersetup{
		pdfborder={0 0 0},
		pdfauthor={\studentname},
		pdfsubject={\thesisname}, 
		pdfcreator={\studentname},
		pdfproducer={\studentname},
	}
}

%Carga el simbolo de grado y el de Angstrom
\newcommand{\angstrom}{\textup{\AA}}
\newcommand{\grad}{$^{\circ}$}
\bibliography{Referencias.bib}
%Inicio del documento, no olvide la etiqueta de cierre al final \end{document}
\begin{document}

%Nombres y formatos de títulos, tablas y figuras
%Use \sffamily para dejar con letra Sans Serif, sin etiqueta queda LaTeX clásico
\renewcommand{\listfigurename}{\sffamily Lista de figuras}
\renewcommand{\listtablename}{\sffamily Lista de tablas}
\renewcommand{\contentsname}{\sffamily Contenido}
\renewcommand{\chaptername}{\sffamily Capítulo}
\renewcommand{\tablename}{\scriptsize \centering \textbf{Tabla}}
\renewcommand{\figurename}{\scriptsize \centering \textbf{Figura}}
\renewcommand{\appendixname}{\sffamily Anexo}

%Cambia el nombre de la sección de referencias
\renewcommand{\bibname}{\sffamily Referencias Bibliográficas}


%Páginas de Presentación del documento - No modificar esto se hace automáticamente
% Páginas de Presentación del documento - Portada mejorada y ajustada
{\newpage
\thispagestyle{empty}
\begin{center}
\begin{figure}
\centering
\epsfig{file=00Figuras/00f00EscudoUN2016.jpg,scale=0.3}% Ajusta la escala del escudo
\end{figure}
\vspace{1cm} % Espacio después del escudo

% Línea decorativa superior
%\rule{\textwidth}{0.8pt} \\[0.3cm]

% Título de la tesis
{\large \textbf{\thesisname}} \\[3.0cm]

% Línea decorativa intermedia
%\rule{\textwidth}{0.5pt} \\[1cm]

% Nombre del estudiante
{\large \textbf{\studentname}} \\[2cm]

% Tipo de documento
{\large \textbf{Proyecto de Tesis Doctoral, requisito parcial para optar al título de Doctor en Estudios
Ambientales.}} \\[2cm]

% Información del director y codirector
{\large \textbf{\director}} \\[2.5cm]
{\large \textbf{Codirectores}}\\
{\large \textbf{\codirector}} \\
{\large \textbf{\codirectora}} \\[2.5cm]


{\large \textbf{UNIVERSIDAD NACIONAL DE COLOMBIA}} \\[0.2cm]
{\large \faculty} \\[0.1cm]
{\large \department} \\[0.1cm]
{\large \sede} \\[0.1cm]

% Línea decorativa inferior
%\rule{\textwidth}{0.8pt} \\[0.5cm]

% Fecha de emisión
{\large \textbf{Diciembre, \issuedate}} \\[0.5cm]
\end{center}
}
% \newpage
% \thispagestyle{empty}
% \vspace{2.5cm}
% \textbf{\ \thesisname} \\
% \vspace{2.5cm}
% \textbf{\Large \studentname} \\
% \vspace{2.5cm}
% \small Proyecto de Tesis
% {\bfseries \academictitle}\\
% \vspace{2.5cm}
% \director \\
% \vspace{2.5cm}
%Codirector(a): \\
%\codirector \\
%\vspace{2.5cm}
%Línea de investigación: \\
%\researchtopic\\
%Grupo de investigación: \\
%\resgroupone \\
%\resgrouptwo \\
%\vspace{2.0cm}
%Universidad Nacional de Colombia \\
%\faculty \\
%\department \\
%\issuedate

% Dedicatorias
% \newpage
% \thispagestyle{empty}
% \begin{flushright}
% \begin{minipage}{12.5cm}
% \noindent
% \\[10em]
% %Modificar la cita que se quiere agregar
% {\Large Cita 01.}
% \\[3em]
% Autor
% \\ \textit{Fuente}
% \\[10em]
% %Para anular la adición de una segunda cita anule las siguientes lineas desde acá mediante comentario (%)
% {\Large \textit{Wenn du es nicht einfach erkl\"{a}ren kannst, hast du es nicht genug verstanden} - Si no eres capaz de explicar algo claramente, es que aún no lo has entendido lo suficiente.}
% \\[3em]
% Albert Einstein
% %Hasta acá!
% \end{minipage}
% \end{flushright}

% Declaracíon de originalidad del texto y del contenido
% No modificar, se hace automáticamente con los comandos ya definidos
% \newpage
% \chapter*{\sffamily Declaración}
% \par Me permito afirmar que he realizado ésta tesis de manera autónoma y con la única ayuda de los medios permitidos y no diferentes a los mencionados el presente texto. Todos los pasajes que se han tomado de manera textual o figurativa de textos publicados y no publicados, los he reconocido en el presente trabajo. Ninguna parte del presente trabajo se ha empleado en ningún otro tipo de tesis.
% \\[1em]
% \sede., \submissiondate
% \\[6em]
% \rule{6cm}{0.5pt}\\
% \studentname
% }

%Páginas preámbulo, listado de figuras, tablas y tabla de contenido
%{\pagestyle{plain} \pagenumbering{roman}
%\setlength{\parskip}{1mm}
%\include{00Agradecimientos}
%\include{00Abreviaturas}
%\newpage
\chapter*{\sffamily Resumen}
\addcontentsline{toc}{chapter}{Resumen}%
\textbf{\Huge \ Título en español } \\
\par Texto del resumen.
\\[2cm]
\textbf{Palabras clave:} \palabrasclave

\newpage 
\chapter*{\sffamily Abstract}
\addcontentsline{toc}{chapter}{Abstract}%
\textbf{\Huge \ Título en inglés} \\
\par Abstract text.
\\[2cm]
\textbf{Keywords:} \keywords

%\newpage 
%\chapter*{\sffamily Zusammenfassung}
%\addcontentsline{toc}{chapter}{Zusammenfassung}%
%\par Zusammenfassung texte.
%\par 
%\\[2cm]
%\textbf{Schlüsselwörter:} \schlusselworter
%\listoffigures
%\addcontentsline{toc}{chapter}{Lista de figuras}
%\listoftables
%\addcontentsline{toc}{chapter}{Lista de tablas}
%\tableofcontents
%\addcontentsline{toc}{chapter}{Contenido}
%\clearpage
%}

\pagenumbering{arabic}
\setlength{\parskip}{\baselineskip}
%Incluir secciones del documento de aqui en adelante
%Use \include para incluir desde una página nueva e \input para incluir sin salto de página
% Archivos incluidos
\chapter{Introducción}

El cambio climático, principalmente resultado de las actividades humanas, se ha convertido en uno de los retos más urgentes de nuestra era \parencite{allanSummaryPolicymakers2023}. Desde la época preindustrial, las concentraciones de gases de efecto invernadero (GEI) como el dióxido de carbono (CO\textsubscript{2}), el metano (CH\textsubscript{4}) y el óxido nitroso (N\textsubscript{2}O) han alcanzado niveles sin precedentes. En 2024, estas concentraciones llegaron a 423.9 partes por millón (ppm) para el CO\textsubscript{2}, 1942 partes por billón (ppb) para el CH\textsubscript{4} y 338.0 ppb para el N\textsubscript{2}O, según las mediciones globales más recientes \parencite{wmoWMOGreenhouseGas2025}. Si bien una concentración atmosférica de CO\textsubscript{2} cercana al 0.042\% puede parecer nominal en términos volumétricos, su relevancia no reside en la abundancia relativa, sino en las propiedades radiativas intrínsecas de la molécula, la cual absorbe y reemite eficazmente la radiación infrarroja térmica. En consecuencia, este incremento acelerado ha alterado el forzamiento radiativo y ha intensificado el calentamiento global, provocando eventos climáticos extremos, alteraciones significativas en los ecosistemas y una rápida pérdida de biodiversidad. 

Dentro de las fuentes antropogénicas de GEI, la gestión inadecuada de residuos orgánicos se destaca como un factor crítico, ya que su descomposición anaeróbica genera grandes cantidades de CH\textsubscript{4}, un gas con un potencial de calentamiento global 28 veces superior al del CO\textsubscript{2} en un horizonte de 100 años. Aunque los océanos y ecosistemas terrestres han actuado históricamente como sumideros naturales, absorbiendo gran parte del CO\textsubscript{2} antropogénico, reportes recientes advierten sobre una reducción en la eficiencia de estos sumideros durante 2024, lo que podría señalar una peligrosa retroalimentación climática. Por tanto, ante los signos de saturación de los mecanismos naturales de absorción y el alto potencial de calentamiento del metano, es urgente desarrollar soluciones innovadoras y sostenibles para mitigar las emisiones actuales, particularmente aquellas vinculadas a la gestión de residuos orgánicos \parencite{ginerFoodLossWaste2025a}.

Las estrategias convencionales para la gestión de residuos orgánicos —como los rellenos sanitarios, la digestión anaerobia y el compostaje— enfrentan limitaciones significativas. Los rellenos sanitarios generan emisiones prolongadas y difícilmente controlables de metano durante años e incluso décadas; la digestión anaerobia, aunque permite capturar biogás, demanda infraestructura compleja y costosa; y el compostaje, a pesar de su mayor accesibilidad, conlleva emisiones de N\textsubscript{2}O y requiere períodos prolongados para lograr la estabilización del material \parencite{kongEvaluatingGreenhouseGas2012,nordahlGreenhouseGasAir2023}. En contextos agroindustriales de países en desarrollo —donde se concentra la mayor parte de los residuos orgánicos sin tratar—, estas alternativas suelen resultar inviables desde el punto de vista técnico o económico. Justamente esta brecha abre una ventana de oportunidad para tecnologías alternativas basadas en procesos biológicos que puedan ser controlados, escalables y adaptados a estas realidades.

Entre las alternativas emergentes, la bioconversión de residuos orgánicos mediante larvas de \textit{Hermetia illucens} (conocida como mosca soldado negra o BSF, por sus siglas en inglés) ha demostrado ser una tecnología prometedora y transformadora. A diferencia de los métodos convencionales, la BSF opera bajo un enfoque de economía circular: transforma residuos orgánicos en biomasa larval rica en proteínas y lípidos y en fertilizante orgánico de alta calidad (\textit{frass}), todo mediante procesos biológicos susceptibles de diseño \parencite{bermudezComprehensiveUtilizationBlack2023,Surendra2016}. Comparada con métodos tradicionales, la bioconversión con BSF presenta ventajas claras: requiere menor infraestructura, opera en tiempos más cortos (3--4 semanas), funciona sin generación de lixiviados contaminantes y genera productos valiosos. Sin embargo, este potencial contrasta con un vacío de información crítica: la mayoría de las implementaciones operan como una ``caja negra'', sin monitoreo ambiental riguroso y sin una cuantificación fiable de los flujos de carbono, lo que impide validar su sostenibilidad real frente al compostaje convencional \parencite{jenkinsProcessingPoultryManure2025}.

Teniendo en cuenta estas limitaciones, se propone la integración de \textbf{tres componentes tecnológicos complementarios}: (1) sensores IoT de bajo costo, (2) modelos matemáticos (EDO) y (3) algoritmos de aprendizaje automático (AAA). Esta convergencia tecnológica no es aditiva, sino \textbf{estructuralmente necesaria} para superar las barreras de cada método por separado: mientras los sensores capturan la realidad física pero carecen de contexto ("ceguera semántica"), y y los modelos matemáticos mecanicistas (caja blanca) aportan coherencia teórica y respeto por las leyes termodinámicas, pero pierden precisión predictiva ante la estocasticidad del comportamiento larval y la heterogeneidad del sustrato sólido, la inteligencia artificial permite cerrar la brecha adaptándose a la estocasticidad del sistema. Es crucial entender que ninguno de estos componentes por sí solo puede validar la sostenibilidad ambiental del proceso; solo su integración híbrida permite transformar datos crudos en evidencia verificable para la toma de decisiones.

El objetivo de esta tesis es desarrollar un sistema híbrido que integre sensores IoT, modelos matemáticos basados en EDO y algoritmos de aprendizaje automático para estimar en tiempo real las emisiones de CO\textsubscript{2} y CH\textsubscript{4} durante la bioconversión de residuos con BSF. El sistema debe: (1) caracterizar dinámicamente las emisiones durante los ciclos productivos; (2) validar predicciones con tolerancia de error aceptable ($<10\%$ MAPE) frente a métodos estándar (cromatografía); y (3) generar herramientas accesibles (dashboard web) que traduzcan datos complejos en información para la toma de decisiones operativas, permitiendo comparar el desempeño ambiental contra una línea base.

El sistema mide: (1) concentraciones de CO\textsubscript{2} y CH\textsubscript{4} mediante sensores NDIR; y (2) variables de proceso (temperatura, humedad, densidad larval estimada). \textbf{Nota sobre el alcance:} No se mide directamente N\textsubscript{2}O en esta etapa por limitaciones de costos en sensores de bajo costo, aunque la arquitectura modular del sistema está diseñada para permitir su integración futura, reconociendo su importancia en el balance global de GEI.

La propuesta, sin embargo, no ignora las complejidades del entorno real. Reconoce la necesidad de gestionar la deriva de los sensores en condiciones de humedad extrema para evitar lecturas erróneas, así como la urgencia de diseñar soluciones robustas ante la falta de datos masivos y la conectividad rural inestable. Además, el uso de inteligencia artificial se somete a un escrutinio ético-ambiental, evaluando si el costo energético del procesamiento y los riesgos de sesgo justifican el beneficio operativo, un balance ineludible en el campo de los Estudios Ambientales.

\section*{Claridad de propósitos: alcance técnico, marco comparativo y responsabilidad contextual}

Esta tesis tiene un alcance técnicamente delimitado pero socialmente situado: desarrollar un sistema híbrido para estimar en tiempo real las emisiones sin pretender sustituir el juicio técnico o normativo. El sistema es, por diseño, una \textit{herramienta descriptiva de soporte a la decisión}, cuyo propósito es proporcionar a productores información cuantitativa para correlacionar decisiones operativas (densidad, sustrato, aireación) con la huella de carbono resultante. El diseño prioriza la accesibilidad y la transparencia, permitiendo generar datos comparables con alternativas como la digestión anaerobia o los rellenos sanitarios.

Reconociendo que la tecnología no es neutral, esta investigación asume una responsabilidad epistemológica explícita. Los algoritmos reflejan los contextos donde se entrenan y rara vez capturan la variabilidad total de los sistemas locales. Por ello, se establecen tres compromisos: (1) \textbf{transparencia algorítmica}, documentando limitaciones y supuestos (IA Explicable); (2) \textbf{soberanía de datos}, permitiendo el acceso y validación por parte de los usuarios locales; y (3) \textbf{compatibilidad con saberes}, asegurando que la tecnología complemente el conocimiento empírico del manejo de residuos. El éxito del sistema no se mide únicamente por métricas de error estadístico, sino por su capacidad para integrarse de manera justa en sistemas agroindustriales reales, fortaleciendo la autonomía técnica del productor.

\subsection*{Aporte del sistema a la gestión ambiental}

Desde la perspectiva de la gestión ambiental, este sistema representa un cambio de paradigma al transitar de tratamientos de residuos basados en la estimación empírica a procesos gobernados por la precisión de los datos. Su principal aporte radica la capacidad de visibilizar el metabolismo del proceso en tiempo real, permitiendo identificar y mitigar picos de emisión de GEI que, en métodos convencionales, pasarían inadvertidos. Al integrar la detección temprana de condiciones anaeróbicas con modelos predictivos, la herramienta no solo optimiza la huella de carbono del tratamiento, sino que otorga trazabilidad y transparencia al ciclo de valorización de residuos, elementos indispensables para validar estrategias de economía circular y cumplimiento normativo en escenarios de crisis climática.
\chapter{Marco Teórico}

El desarrollo de soluciones tecnológicas para la gestión sostenible de residuos orgánicos exige una comprensión integrada de los procesos biológicos, ambientales y computacionales involucrados. Este capítulo ofrece la base teórica que sustenta la propuesta de un sistema inteligente híbrido para la predicción de emisiones de gases de efecto invernadero (GEI) en la bioconversión con BSF. Se articulan dos ejes interdependientes: (i) la bioconversión como estrategia de valorización de residuos en contextos agroindustriales; (ii) los enfoques de modelado matemático y sistemas inteligentes para la estimación dinámica de emisiones.

\section{Residuos orgánicos y bioconversión con \textit{Hermetia illucens}}

En contextos agroindustriales del sur global —como el Valle del Cauca—, los residuos orgánicos provenientes de cultivos como café, caña y frutas se generan de forma dispersa, en picos estacionales y con alta carga biodegradable, lo que incrementa el riesgo de emisiones de CH\textsubscript{4} cuando su gestión es inadecuada \parencite{FAO2021,Leddin2024}. Frente a las limitaciones de las estrategias convencionales —rellenos sanitarios, digestión anaerobia y compostaje—, que suelen ser inviables técnica o económicamente en estos entornos \parencite{kongEvaluatingGreenhouseGas2012,nordahlGreenhouseGasAir2023}, la bioconversión con BSF se ha adoptado crecientemente como una alternativa de economía circular.

La bioconversión de residuos orgánicos mediante larvas en estadios voraces (L3--L5) transforma desechos en biomasa rica en proteínas y lípidos, así como en \textit{frass}, un fertilizante orgánico de alto valor agronómico \parencite{Surendra2016,bermudezComprehensiveUtilizationBlack2023,singhInclusiveApproachOrganic2019}. Este proceso se sustenta en una sinergia entre la acción macroscópica de las larvas y la actividad microbiana del sustrato: las larvas fragmentan mecánicamente el material, aumentando su área superficial, mientras su microbioma intestinal ---junto con microorganismos del medio--- degrada enzimáticamente polímeros complejos en compuestos asimilables. Esta doble acción reduce entre el 50\,\% y el 80\,\% la masa del residuo (base húmeda) y altera drásticamente las propiedades fisicoquímicas del medio (pH, humedad, disponibilidad de nutrientes), generando microambientes dinámicos que condicionan si predominan rutas metabólicas aeróbicas o anaeróbicas. Durante este proceso, el dióxido de carbono (\ce{CO2}) se genera principalmente por la respiración aeróbica tanto de las larvas como de la microbiota asociada; estudios estiman que hasta un tercio del carbono ingerido se libera como \ce{CO2} \parencite{dienerConversionOrganicMaterial2009}, una emisión altamente sensible a factores como la composición del sustrato, densidad larval, temperatura y oxigenación \parencite{rossiEstimatingDynamicsGreenhouse2024,bekkerImpactSubstrateMoisture2021a,eriksenDynamicModellingFeed2022}.
%\begin{figure}
%    \centering
%    \includegraphics[scale=0.15]{clico-mosca.jpg}
%    \caption{Ciclo de vida de la mosca soldado negra (\textit{Hermetia illucens}), mostrando sus estadios de desarrollo.}
%    \label{fig:ciclo-mosca}
%\end{figure}

Aunque la bioconversión muestra un perfil de emisiones potencialmente más favorable que el compostaje tradicional ---con reducciones reportadas en \ce{CH4} y \ce{N2O} \parencite{jenkinsProcessingPoultryManure2025}---, su implementación a escala carece de sistemas rigurosos para cuantificar y predecir en tiempo real las emisiones gaseosas asociadas. Esta limitación impide su validación como solución climática verificable y obstaculiza su integración en mecanismos de compensación de carbono o políticas de gestión sostenible de residuos. En particular, la dinámica de emisiones, especialmente la de \ce{CO2} ---que refleja directamente la actividad biológica en curso---, no ha sido suficientemente modelada mediante herramientas matemáticas capaces de ofrecer estimaciones continuas y predictivas bajo condiciones operativas variables. Así, persiste una brecha crítica entre el potencial ambiental demostrado en ensayos controlados y la capacidad de monitoreo, verificación y escalabilidad requerida para su reconocimiento en estrategias climáticas basadas en evidencia.

Aunque el metano (\ce{CH4}) es un gas de efecto invernadero con un potencial de calentamiento global significativamente mayor que el del dióxido de carbono (\ce{CO2}), su generación en sistemas de bioconversión con BSF es generalmente baja debido al carácter predominantemente aeróbico del proceso. No obstante, se incluye su monitoreo como variable de control para detectar condiciones transitoriamente anaeróbicas ---por ejemplo, por sobresaturación del sustrato o deficiente aireación--- que podrían comprometer la eficiencia del proceso y generar picos no deseados de emisiones. Sin embargo, el foco técnico de esta tesis recae en el \ce{CO2}, ya que constituye el principal indicador cuantitativo y dinámico de la actividad biológica total (tanto larval como microbiana) durante la bioconversión. Además, presenta una señal más estable y continua, lo que facilita su integración en modelos de estimación en tiempo real, y permite una correlación directa con variables operativas clave, como el consumo de sustrato, la biomasa generada y las tasas respiratorias. Por estas razones, el \ce{CO2} se adopta como la variable central para el desarrollo, calibración y validación del sistema híbrido propuesto.

La aptitud de los residuos para la bioconversión varía significativamente según su composición. Como se resume en la Tabla~\ref{tab:residuos_BSF}, subproductos agroindustriales como la pulpa de café o cáscaras de frutas son altamente adecuados, mientras que residuos lignocelulósicos o lodos orgánicos requieren pretratamiento o presentan riesgos sanitarios \parencite{elsayedConversionProteinrichWaste2024,brunoValorizationOrganicWaste2025}. Esta heterogeneidad —común en el Valle del Cauca— subraya la necesidad de sistemas de monitoreo adaptables, capaces de caracterizar el proceso en función del sustrato y las condiciones locales.

La evidencia cuantitativa disponible es fragmentada: se reportan emisiones de CO\textsubscript{2} entre 7.76 y 11.88\,g por gramo de larva seca \parencite{rossiEstimatingDynamicsGreenhouse2024}, o 96.5\,g de CO\textsubscript{2} por kg de residuo tratado \parencite{fuhrmannComprehensiveIndustryrelevantBlack2025}, pero sin protocolos estandarizados ni mediciones continuas. Esta incertidumbre impide comparar tecnologías, diseñar políticas públicas informadas o acceder a mecanismos de mercado basados en la mitigación climática.

\begin{table}[H]
\centering
\caption{Clasificación de residuos orgánicos y su aptitud para bioconversión con \textit{Hermetia illucens}}
\label{tab:residuos_BSF}
\begin{tabular}{|p{4cm}|p{5cm}|p{5cm}|}
\hline
\textbf{Tipo de residuo} & \textbf{Origen común} & \textbf{Aptitud para \textit{BSF}} \\ \hline
Residuos agrícolas & Restos de cosecha, podas, hojas secas & Moderada. Requieren pretratamiento por bajo contenido proteico. \\ \hline
Residuos alimentarios & Sobras de cocina, frutas, verduras, pan & Alta. Contenido balanceado de lípidos y carbohidratos. \\ \hline
Subproductos agroindustriales & Pulpa de café, cáscaras de frutas, bagazo de caña & Alta. Composición variable, pero generalmente adecuada. \\ \hline
Grasas y aceites residuales & Aceites usados, residuos grasos de cocina & Media. Requieren mezcla con otros residuos para evitar inhibición. \\ \hline
Lodos orgánicos & Lodos de depuradoras, biosólidos & Baja. Riesgo sanitario y composición inconsistente. \\ \hline
\end{tabular}
\end{table}

\textbf{Hacia sistemas de monitoreo adaptable en contextos heterogéneos:}

La heterogeneidad de residuos agroindustriales documentada en la Tabla~\ref{tab:residuos_BSF} 
—desde pulpa de café (alta aptitud) hasta lodos orgánicos (baja aptitud)— refleja una realidad 
común en el Valle del Cauca: los productores raramente procesan un único tipo de residuo de forma 
consistente. En su lugar, enfrentan mezclas variables según la estacionalidad, disponibilidad local 
y dinámicas de producción \parencite{elsayedConversionProteinrichWaste2024,brunoValorizationOrganicWaste2025}. 
Adicionalmente, variables ambientales críticas —temperatura, humedad del sustrato, oxigenación— 
fluctúan durante el ciclo de bioconversión, particularmente en contextos rurales donde el control 
es limitado \parencite{bekkerImpactSubstrateMoisture2021a,rossiEstimatingDynamicsGreenhouse2024}. 

Esta variabilidad operativa implica que los métodos de medición estática o puntual —como 
muestreos puntuales con cromatografía de gases— son insuficientes para caracterizar emisiones 
de manera confiable. Se requieren sistemas de monitoreo que sean: (i) \textbf{continuos en tiempo 
real}, capaces de capturar dinámicas horarias y diarias; (ii) \textbf{adaptables a diferentes tipos 
de residuos}, disminuyendo la necesidad de recalibración manual frecuentemente; (iii) \textbf{interpretables}, 
permitiendo a productores locales comprender qué sucede en el sistema y tomar decisiones operativas 
informadas. Es precisamente esta necesidad la que motiva un enfoque integrado que combine 
modelación matemática dinámica, sensores de bajo costo y calibración automática mediante algoritmos 
de aprendizaje automático, tema de la siguiente sección.

\section{Modelación de emisiones y sistemas inteligentes}

La cuantificación dinámica de emisiones en procesos biológicos exige herramientas que articulen una base mecanística con capacidad de adaptación empírica. En este contexto, el modelado matemático basado en ecuaciones diferenciales ordinarias (EDO) ha impulsado avances significativos. Por ejemplo, el modelo de crecimiento larval propuesto por Eriksen (2022) integra balances energéticos y cinética logística para simular simultáneamente la acumulación de biomasa y la producción de \ce{CO2}, combinando los principios del enfoque \textit{Dynamic Energy Budget} (DEB) con eventos biológicos clave del ciclo larval, como las mudas y la metamorfosis. De manera análoga, funciones cinéticas de tipo Monod y Haldane —originalmente desarrolladas para digestión anaerobia— han sido adaptadas con éxito para describir la degradación del sustrato y el crecimiento microbiano en sistemas de bioconversión con insectos \parencite{winAnaerobicDigestionBlack2018, ahlamineMathematicalAnalysisAnaerobic2024}. 

Si bien los modelos mecanicistas actuales proporcionan una base teórica indispensable, su aplicación directa en entornos operativos dinámicos puede presentar limitaciones de alcance. La literatura \parencite{bekkerImpactSubstrateMoisture2021a,rossiEstimatingDynamicsGreenhouse2024} sugiere que, al validarse principalmente bajo condiciones controladas, las formulaciones clásicas a menudo requieren simplificaciones sobre las interacciones simultáneas de variables (temperatura, humedad, oxigenación). Asimismo, el uso de parámetros fijos podría restringir la capacidad del modelo para responder a la variabilidad estocástica del sustrato en campo. Esta diferencia entre las condiciones controladas y la operación real sugiere la conveniencia de explorar enfoques híbridos o basados en datos que complementen la teoría existente con información en tiempo real.

Paralelamente, los avances en sensores de bajo costo y plataformas IoT han 
facilitado el monitoreo continuo de variables críticas —CO\textsubscript{2}, 
CH\textsubscript{4}, temperatura, humedad— en entornos operativos reales 
\parencite{rajRealTimeEstimation2023,duobieneDevelopmentWirelessSensor2022,
yavariArtEMonArtificialIntelligence2023}. Estos datos, procesados mediante 
algoritmos de aprendizaje automático (AAA), permiten identificar patrones 
complejos y no lineales que serían difíciles de capturar mediante reglas 
explícitas \parencite{weichertReviewMachineLearning2019,chenOptimizationModelProcess2021}. 

En particular, la integración de modelos mecanísticos (EDO) con AAA, enfoque conocido en la literatura como \textbf{modelado de caja gris, (grey box modelling)} ha demostrado 
potencial en sistemas biotecnológicos complejos: mientras el modelo EDO proporciona 
una estructura interpretable y basada en principios físicos, el AAA permite calibración 
dinámica de parámetros a partir de datos en tiempo real, mejorando predicciones cuando 
las condiciones operativas se desvían del escenario de entrenamiento 
\parencite{kobelskiModelbasedProcessOptimization2024a,guillaumeAsymptoticEstimatedDigestibility2023}. 

Sin embargo, esta integración enfrenta desafíos prácticos: (i) los algoritmos de AAA 
requieren suficientes datos históricos para entrenar sin riesgo de \textit{overfitting}, 
pero ciclos de bioconversión típicos generan pocas muestras por evento; (ii) la validación 
cruzada y selección de hiperparámetros requieren cuidado especial en contextos con datos 
limitados; (iii) la elección entre arquitecturas simples (regresión lineal múltiple) versus 
complejas (redes neuronales) dependerá del volumen y calidad de datos disponibles. Por ello, 
esta tesis adopta un enfoque escalonado: comenzar con modelos interpretables (regresión múltiple), 
evaluar redes neuronales multicapa solo si se dispone de base de datos amplia, y validar todas 
las predicciones contra mediciones de referencia de laboratorio.

En síntesis, el marco teórico desarrollado en este capítulo fundamenta un enfoque 
técnico híbrido que integra tres componentes interdependientes: (i) un modelo mecanístico 
basado en ecuaciones diferenciales ordinarias que captura la dinámica biológica del proceso 
de bioconversión; (ii) sensores IoT de bajo costo que permiten monitoreo continuo de emisiones 
de CO\textsubscript{2} y CH\textsubscript{4} en tiempo real; (iii) algoritmos de aprendizaje 
automático que facilitan calibración dinámica y predicción adaptativa bajo variaciones 
operativas. Esta integración es particularmente relevante para contextos agroindustriales 
del sur global, donde la falta de infraestructura costosa y la necesidad de soluciones 
escalables hacen que un sistema de bajo costo y mantenible sea más viable que tecnologías 
convencionales.

Operacionalmente, la ``inteligencia'' del sistema reside en su capacidad para: (1) generar 
predicciones con fundamento mecanístico (interpretables, no cajas negras); (2) actualizar 
predicciones automáticamente a partir de datos en tiempo real; (3) proporcionar información 
clara a productores locales para que tomen decisiones informadas sobre condiciones operativas 
(temperatura, humedad, manejo de residuos). No obstante, esta investigación reconoce 
explícitamente que la mera disponibilidad de datos técnicos no garantiza equidad o sostenibilidad. 
La verdadera gobernanza de un sistema como este requiere —más allá del alcance de esta tesis— 
decisiones sobre acceso a datos, beneficios compartidos, y co-diseño con usuarios finales. 
Estas dimensiones ético-políticas constituyen un campo de investigación complementario 
igualmente importante.

\section{Gobernanza de datos, ética algorítmica y epistemología ambiental}

La implementación de tecnologías digitales en la gestión de procesos biológicos no es un acto neutral; constituye la creación de un sistema socio-técnico donde convergen biología, código y decisión humana. Desde la perspectiva de los estudios de ciencia, tecnología y sociedad (CTS) y la epistemología ambiental, es imperativo trascender la visión instrumental de la tecnología para abordar las implicaciones éticas y cognitivas de ``datificar'' la naturaleza.

\subsection{De la opacidad algorítmica a la interpretabilidad basada en modelos}
Uno de los mayores riesgos en la aplicación de la inteligencia artificial (IA) al medio ambiente es la \textit{opacidad algorítmica} \parencite{burrelHowMachineThinks2016}. En modelos de caja negra, las decisiones se vuelven inescrutables, generando relaciones de dependencia tecnológica donde el operario obedece sin comprender.

Esta tesis adopta una postura ética fundada en la \textit{interpretabilidad por diseño}. Al integrar modelos mecanísticos basados en ecuaciones diferenciales ordinarias (EDO) —cuyos parámetros tienen significado biológico explícito— con algoritmos de aprendizaje ligeros, el sistema no solo predice eventos (el \textit{qué}), sino que permite rastrear sus causas fisiológicas o operativas (el \textit{porqué}), como el exceso de humedad o la acumulación de zonas anaeróbicas. Esta transparencia cognitiva devuelve la agencia al tomador de decisiones y se alinea con los principios de transparencia exigidos por marcos de reporte climático como el IPCC Tier 3 \parencite{adadiPeekingBlackBoxSurvey2018}.

\subsection{Soberanía de datos y justicia algorítmica en el Sur Global}
La gobernanza de los datos ambientales plantea interrogantes críticos sobre propiedad, acceso y beneficio. La literatura sobre \textit{Big Data} en agricultura alerta sobre el ``extractivismo de datos'', donde la información generada en territorios rurales es capturada por plataformas externas sin retorno de valor local \parencite{carolanAutomatedAgrifoodFutures2020}.

Esta investigación se alinea con los principios de \textit{soberanía de datos} y \textit{diseño frugal}. La arquitectura propuesta es descentralizada: el procesamiento crítico ocurre en el borde (\textit{edge computing}), minimizando la dependencia de infraestructuras en la nube centralizadas y costosas. Esto no solo reduce costos, sino que garantiza que los datos sensibles permanezcan bajo el control del productor. Filosóficamente, el monitoreo ambiental no debe ser una herramienta de vigilancia, sino un instrumento de empoderamiento que permita a los productores agroindustriales validar sus prácticas ante mercados globales, sin alienar su conocimiento local.

\subsection{Límites del reduccionismo tecnológico}
Finalmente, se reconoce la tensión ontológica entre la complejidad del sistema biológico y su representación digital. Un sensor de \ce{CO2} registra una señal fisicoquímica, pero no captura dimensiones cualitativas como la resiliencia del microbioma o el bienestar larval. Existe el riesgo de caer en el \textit{solucionismo tecnológico}, asumiendo que lo no medido es irrelevante.

Por ello, este marco teórico establece que el sistema híbrido es una herramienta de \textit{soporte}, no de \textit{sustitución}. La comprensión más robusta del proceso emerge de la interacción dialógica entre la señal del sensor y la experiencia empírica del operador. El modelo matemático y la inteligencia artificial actúan como traductores que hacen visible lo invisible (los gases), pero la interpretación final debe permanecer anclada en el contexto ecológico, social y productivo del territorio.
\chapter{Pregunta de Investigación}

¿En qué medida la integración de un modelo matemático basado en ecuaciones diferenciales ordinarias, sensores IoT de bajo costo y algoritmos de aprendizaje automático permite desarrollar un sistema híbrido que no solo estime con precisión las emisiones de \ce{CO2} y \ce{CH4}, sino que constituya una herramienta de monitoreo dinámico para fundamentar la toma de decisiones operativas y generar métricas verificables del desempeño ambiental en la bioconversión de residuos agroindustriales con \textit{Hermetia illucens}?

\chapter{Hipótesis de Investigación}

\section{Hipótesis General}

Un sistema híbrido que integre modelamiento matemático basado en ecuaciones diferenciales ordinarias (EDO), sensores IoT de bajo costo y algoritmos de aprendizaje automático permite estimar las emisiones de \ce{CO2} y \ce{CH4} en tiempo real con un error absoluto medio porcentual (MAPE) inferior al 10\,\% respecto a mediciones de referencia. Esta capacidad de monitoreo continuo proporciona, a su vez, indicadores dinámicos sobre la eficiencia metabólica del proceso, permitiendo identificar ineficiencias operativas (como la anaerobiosis) y cuantificar el desempeño ambiental del sistema más allá de las estimaciones estáticas tradicionales.

\section{Hipótesis Específicas}
\subsection{H1: Precisión instrumental de bajo costo}
Los sensores NDIR de bajo costo, sometidos a protocolos de calibración y preacondicionamiento de señal, permiten medir concentraciones de \ce{CO2} y \ce{CH4} en el entorno de bioconversión con una desviación (MAPE) $\leq 10$\,\% respecto a equipos certificados de laboratorio.

\subsection{H2: Capacidad predictiva del modelo mecanístico}
El modelo de EDO, fundamentado en la bioenergética del crecimiento larval y parametrizado con datos locales, es capaz de predecir la tasa de producción de \ce{CO2} durante el ciclo experimental de 15 días con un error porcentual absoluto medio (MAPE) $\leq 15$,\%. Esta validación, realizada en ventanas de observación de 30 minutos cada dos días, demuestra que el modelo captura correctamente la dinámica metabólica subyacente del sistema.

\subsection{H3: Refinamiento por aprendizaje automático}
La integración de una capa de corrección mediante algoritmos de aprendizaje automático reduce el error residual del modelo de EDO en al menos un 30\,\%, logrando que el sistema híbrido alcance un MAPE $< 10$\,\% ante perturbaciones no modeladas.
\subsection{H4: Valor para la gestión operativa} La incorporación de medición directa de \ce{CH4} en el esquema de monitoreo permite la detección inmediata de eventos transitorios de anaerobiosis (picos de emisión) que los modelos estáticos basados en factores de emisión promedio omiten. Esto proporciona indicadores dinámicos que habilitan acciones correctivas en tiempo real, reduciendo la incertidumbre operativa frente a las estimaciones puramente teóricas.

\emph{Nota: Las métricas de precisión (H1--H3) se evaluarán bajo condiciones experimentales de laboratorio (30\,°C $\pm$ 2\,°C, HR 60--70\,\%), mientras que la H4 se validará mediante análisis de escenarios operativos.}
\chapter{Planteamiento del Problema}

En el contexto Tecnológico y Biológico, la gestión de la fracción orgánica de los residuos sólidos municipales (FORSU) enfrenta limitaciones estructurales en sus rutas convencionales. Los rellenos sanitarios operan como reactores anaerobios no controlados de larga duración; el compostaje, aunque accesible, requiere extensos tiempos de estabilización y conlleva riesgos de emisiones de óxido nitroso (N\textsubscript{2}O); y la digestión anaerobia, pese a su ventaja energética, demanda inversiones de capital, a menudo inviables para la escala agroindustrial media.

En este escenario, la bioconversión mediante larvas de \textit{Hermetia illucens} (BSF) se ha consolidado como una alternativa de alta eficiencia. A diferencia de los métodos pasivos, este es un sistema biológico acelerado: durante su fase larval (aprox. 14 días), el insecto incrementa su masa corporal hasta 4000 veces mediante un metabolismo voraz y exotérmico que ocurre en simbiosis con una microbiota activa. Este proceso transforma la materia orgánica en biomasa (proteína) y fertilizante (\textit{frass}) con baja huella hídrica, posicionándose como una herramienta clave para la economía circular \parencite{Surendra2016,bermudezComprehensiveUtilizationBlack2023}.

Sin embargo, a pesar de estas ventajas biológicas, la ausencia de cuantificación dinámica y rigurosa de las emisiones de gases de efecto invernadero (GEI), particularmente CH\textsubscript{4} y CO\textsubscript{2}, genera una incertidumbre crítica. Actualmente, la tecnología opera mayoritariamente como una ``caja negra'': se conocen los insumos y productos, pero se ignora la dinámica de los flujos gaseosos intermedios, basándose en factores de emisión estáticos o extrapolaciones de otros procesos (compostaje) que no reflejan la realidad metabólica de la BSF \parencite{fuhrmannComprehensiveIndustryrelevantBlack2025,rossiEstimatingDynamicsGreenhouse2024}. 

Esta fragmentación del conocimiento tiene consecuencias técnicas, operativas y de gobernanza climática:

\textbf{Primero}, al no existir factores de emisión específicos para BSF en las metodologías del IPCC ni en el GHG Protocol, cualquier balance de carbono se basa en supuestos no verificados. Esto introduce sesgos sistemáticos que pueden subestimar significativamente las emisiones de CH\textsubscript{4}, especialmente bajo condiciones de manejo subóptimas. La consecuencia no es solo técnica, sino ética: se corre el riesgo de asignar beneficios climáticos no reales a una tecnología que, sin monitoreo, podría superar las emisiones del compostaje en escenarios de anaerobiosis no controlada (Greenwashing involuntario).

\textbf{Segundo}, la imposibilidad de detectar en tiempo real picos de CH\textsubscript{4} impide identificar fallas operativas críticas —como sobrepoblación larval, exceso de humedad o deficiente aireación— que generan microambientes anaeróbicos. Estas condiciones no solo aumentan las emisiones de GEI, sino que reducen la eficiencia de conversión, incrementan la mortalidad larval y degradan la calidad del \textit{frass}, impactando directamente la viabilidad económica del sistema. Medir las emisiones es, por tanto, una herramienta de control de proceso.

\textbf{Tercero}, la carencia de sistemas de Medición, Reporte y Verificación (MRV) alineados con estándares internacionales constituye una barrera estructural para el acceso a mecanismos de financiación climática. Según el Banco Mundial, los mercados de carbono modernos exigen trazabilidad cuantitativa y verificable \parencite{worldbankStateTrendsCarbon2023}. Normativas como el \textit{Verified Carbon Standard} (VCS) estipulan que la cuantificación debe ser "real, medible y permanente" \parencite{verraVCSStandardV472024}. Sin datos empíricos robustos, la industria de insectos no puede certificar sus reducciones a escala industrial \parencite{smetanaSustainableUseHermetia2019}.

Esta brecha se ve agravada por la escasa integración de tecnologías de monitoreo con modelos predictivos. La mayoría de las implementaciones actuales carecen de retroalimentación en tiempo real, operando con parámetros fijos que ignoran la dinámica no lineal del proceso biológico \parencite{kamilaris2017review,liuOptimizingDataPipelines2023}.

Frente a este escenario, esta tesis propone un sistema híbrido que transforme la cuantificación de emisiones en una herramienta de gestión ambiental y operativa. Al estimar de forma continua y precisa las emisiones de CO\textsubscript{2} y CH\textsubscript{4}, el sistema permitirá: (1) generar la evidencia empírica necesaria para establecer factores de emisión específicos y alinear la BSF con marcos climáticos globales; (2) optimizar el proceso mediante la detección temprana de ineficiencias a partir de los gases medidos; y (3) construir una trazabilidad verificable que permita hacia el futuro la participación de los productores agroindustriales en economías circulares con credenciales climáticas robustas.
\chapter{Antecedentes}

A continuación, se presentan los principales antecedentes relacionados con el estudio de esta propuesta doctoral, enfocados en la problemática de gestión de residuos, el potencial de soluciones biotecnológicas y, crucialmente, una revisión crítica sobre el papel y las limitaciones de las tecnologías emergentes en la mitigación del impacto ambiental.

\section{Problemática global en el manejo de residuos y la generación de emisiones de GEI}

La gestión inadecuada de residuos constituye una de las principales amenazas ambientales a escala global, con efectos directos sobre el cambio climático, la salud pública y la integridad de los ecosistemas. Según la FAO, se estima que más del 50\,\% de los residuos sólidos generados en países de ingresos medios y bajos son de origen orgánico \parencite{FAO2021}, los cuales, al ser dispuestos sin tratamiento adecuado, se degradan en condiciones anaeróbicas generando metano (\ce{CH4}). Este fenómeno es particularmente grave en sectores como la agricultura, donde las prácticas tradicionales —como la quema de residuos y el uso intensivo de insumos químicos— agravan la carga ambiental de los sistemas productivos.

De acuerdo con el Quinto Informe del Panel Intergubernamental sobre Cambio Climático, aproximadamente el 24\,\% de las emisiones globales de GEI provienen de actividades relacionadas con el uso del suelo, la agricultura y los cambios en el uso de la tierra \parencite{IPCC2014}. A estas emisiones se suman los impactos generados por el manejo ineficiente de residuos en contextos urbanos e industriales, incluyendo lixiviación contaminante, contaminación atmosférica y afectación de la biodiversidad local. Estas condiciones se ven agravadas por la limitada infraestructura de tratamiento y las brechas en gobernanza ambiental, especialmente en regiones rurales y en países en desarrollo.

Este contexto exige de manera urgente un diseño de estrategias sostenibles de gestión de residuos, que integren enfoques sistémicos y soluciones basadas en principios de economía circular. El enfoque de “residuo como recurso” promueve la valorización de la materia orgánica a través de tecnologías como el compostaje, la digestión anaerobia, la vermicultura y, más recientemente, la bioconversión con insectos. Estos enfoques permiten cerrar ciclos de nutrientes, reducir emisiones y generar productos útiles como biofertilizantes, energía o proteínas alternativas, entre otros.

\section{Tecnologías emergentes: Potencialidades y análisis crítico de riesgos}

Para superar las barreras de medición y control en estos bioprocesos, la literatura reciente sugiere la integración de sistemas de gestión inteligentes, capaces de monitorear en tiempo real la generación de emisiones. Esto implica la incorporación de tecnologías como sensores de Internet de las Cosas (IoT), modelos matemáticos y algoritmos de aprendizaje automático, herramientas que prometen mejorar la eficiencia y garantizar la trazabilidad basada en datos \parencite{ahmedAgricultureClimateChange2020}.

Sin embargo, la adopción de estas tecnologías no está exenta de riesgos y limitaciones que deben ser analizados críticamente para evitar soluciones contraproducentes (“tecnosolucionismo”). Diversos autores advierten sobre tres desafíos fundamentales:

\textbf{1. Huella ambiental del hardware:} La implementación masiva de sensores IoT conlleva el riesgo de generar residuos electrónicos (e-waste) debido a la corta vida útil de los componentes en ambientes corrosivos como los de la bioconversión (alta humedad y presencia de amoniaco). Además, la fabricación y operación de estos dispositivos tiene su propia huella de carbono asociada, que debe ser compensada por las eficiencias que el sistema genera \parencite{dharCarbonImpactArtificial2020}.

\textbf{2. Robustez y Sesgo Algorítmico:} Los modelos de aprendizaje automático entrenados en condiciones controladas de laboratorio a menudo fallan al generalizar en entornos rurales reales, donde la variabilidad de los residuos y las condiciones climáticas es alta. Confiar ciegamente en predicciones algorítmicas sin una validación física continua puede llevar a decisiones de manejo erróneas \parencite{klerkxReviewSocialScience2019}.

\textbf{3. Brecha de Implementación:} En contextos del Sur Global, la dependencia de conectividad estable y soporte técnico especializado puede convertirse en una barrera de adopción para pequeños productores. Por tanto, la innovación tecnológica debe equilibrarse con criterios de robustez, bajo costo y mantenibilidad local para ser verdaderamente sostenible.

\section{Estudios Previos sobre la Gestión de Residuos con \textit{BSF}}

En los últimos años, la literatura científica ha validado el potencial de BSF no solo como agente de reducción de residuos, sino como bioreactor eficiente. Estudios fundamentales, como los de \textcite{salamEffectDifferentEnvironmental2022}, han caracterizado la sensibilidad del proceso a variables abióticas, estableciendo que rangos de temperatura entre $25-30\ ^{\circ}\mathrm{C}$ y humedad del sustrato mayores que $60\%$ son críticos para maximizar la tasa metabólica. \textit{Esta dependencia no lineal entre las condiciones ambientales y el rendimiento biológico subraya la necesidad de sistemas de control dinámico, ya que desviaciones menores en estas variables pueden alterar drásticamente la eficiencia de bioconversión.}

La transición hacia modelos predictivos ya cuenta con precedentes importantes que es necesario analizar. En el ámbito del modelado matemático, en \parencite{eriksenDynamicModellingFeed2022} desarrollaron una aproximación mecanística basada en la teoría de Presupuestos Energéticos Dinámicos (DEB), logrando simular el crecimiento larval en función del alimento disponible. Sin embargo, este modelo se centra exclusivamente en la acumulación de biomasa, dejando de lado la predicción dinámica de los flujos gaseosos asociados. Paralelamente, en el dominio del hardware, autores como \textcite{rajRealTimeEstimation2023} y \textcite{duobieneDevelopmentWirelessSensor2022} han implementado prototipos de monitoreo IoT para variables ambientales básicas (temperatura, humedad). Aunque estos trabajos demuestran la viabilidad técnica de la sensorización, se limitan a la vigilancia pasiva de condiciones, careciendo de una capa de inteligencia artificial que integre la variable biológica con la ambiental para predecir anomalías en tiempo real.

Si bien investigaciones recientes han ampliado el horizonte de valorización hacia biocombustibles y bioplásticos \parencite{koyunogluBiofuelProductionUtilizing2024,suryatiLauricAcidBlack2023}, la mayoría de estos trabajos adopta un enfoque de “caja negra”, analizando entradas y salidas sin modelar la dinámica interna del proceso. Persiste, por tanto, una brecha notable: la falta de integración entre los modelos biológicos profundos (tipo EDO/DEB) y la captura de datos en tiempo real (IoT) para la predicción específica de GEI, lo cual limita la capacidad de intervención preventiva en sistemas industriales.

\section{Medición y Dinámica de Emisiones de GEI en la Bioconversión}

La cuantificación de GEI es el pilar para la validación ambiental de la tecnología. Aunque el proceso es predominantemente aeróbico —generando principalmente \ce{CO2} y calor metabólico—, la heterogeneidad del sustrato puede inducir la formación de microambientes anaeróbicos productores de \ce{CH4} y \ce{N2O}, gases con un potencial de calentamiento global significativamente mayor \parencite{bekkerImpactSubstrateMoisture2021a}.

La evidencia cuantitativa actual es dispersa. Por ejemplo, \textcite{jenkinsProcessingPoultryManure2025} reportan reducciones netas de emisiones frente al compostaje debido a ciclos de tratamiento más cortos. En términos de factores de emisión, \textcite{fuhrmannComprehensiveIndustryrelevantBlack2025} reportan $96.5\,\mathrm{g}$ \ce{CO2} por kg de residuo, mientras que \textcite{rossiEstimatingDynamicsGreenhouse2024} presentan valores normalizados por biomasa seca ($7.76$ -- $11.88\,\mathrm{g}$ \ce{CO2}/g larva). \textit{Esta disparidad en las unidades funcionales y la falta de protocolos estandarizados dificultan la comparación directa y el establecimiento de líneas base para mecanismos de compensación de carbono.}

La mayoría de estos estudios se basan en muestreos puntuales (cámaras estáticas) que no capturan los picos transitorios de emisión. Aquí radica una oportunidad para ésta investigación: la implementación de sistemas de monitoreo continuo (IoT) no solo permitiría refinar estos factores de emisión, sino utilizar el \ce{CO2} como indicador directo de la actividad metabólica y el \ce{CH4} como indicador temprano de fallos operativos (anaerobiosis), transformando la medición ambiental en una herramienta de control de procesos \parencite{kamilaris2017review,zhangApplicationBigData2024}.

\section{Antecedentes en Modelamiento Matemático de Procesos de Bioconversión}

El modelamiento matemático constituye la base teórica para transitar de la observación empírica a la predicción cuantitativa en procesos biológicos. En el contexto de la bioconversión con BSF, los modelos mecanísticos permiten desacoplar los fenómenos de crecimiento larval de la dinámica de degradación del sustrato, estableciendo relaciones causales entre la entrada de materia orgánica y la generación de gases de efecto invernadero (GEI), principalmente \ce{CO2} y, en condiciones subóptimas, \ce{CH4}.

La literatura actual se ha centrado predominantemente en modelos deterministas basados en sistemas de Ecuaciones Diferenciales Ordinarias (EDO). Referentes como \textcite{padmanabhaComprehensiveDynamicGrowth2020} adaptaron el marco de Presupuestos Energéticos Dinámicos (DEB), logrando describir con alta precisión los flujos metabólicos a nivel individual (respiración, asimilación y excreción). Posteriormente, \textcite{eriksenDynamicModellingFeed2022} escalaron estos principios a nivel poblacional, incorporando balances de masa para predecir la producción de calor y \ce{CO2}. Sin embargo, estos modelos asumen condiciones de homogeneidad ideal en el sustrato, lo que limita su capacidad para predecir la formación de microambientes anaeróbicos y, por ende, la generación estocástica de metano.

Para abordar la cinética de degradación, otros autores han extrapolado funciones de tipo Monod y Haldane desde la digestión anaerobia \parencite{winAnaerobicDigestionBlack2018, ahlamineMathematicalAnalysisAnaerobic2024}. Si bien estas ecuaciones capturan las no linealidades del crecimiento microbiano y la inhibición por sustrato, su aplicación directa en matrices sólidas porosas (como el lecho de cría de BSF) introduce errores estructurales al no considerar la variabilidad espacial de la temperatura y la humedad \parencite{rossiEstimatingDynamicsGreenhouse2024}.

En la tesis, los modelos mecanísticos actuales (EDO/DEB) son excelentes para describir el comportamiento “promedio” o ideal del sistema (tendencia base), pero carecen de la flexibilidad plástica para adaptarse a perturbaciones no modeladas o variaciones abruptas en la composición del residuo. La literatura evidencia una ausencia de arquitecturas híbridas que utilicen datos de sensores en tiempo real para recalibrar dinámicamente estos parámetros cinéticos, una estrategia necesaria para garantizar la precisión predictiva en escenarios de gestión operativa real.

\section{Avances en IoT y AAA para la gestión sostenible de residuos}

La incorporación de tecnologías basadas en IoT y AAA ha demostrado tener un impacto tangible en la sostenibilidad de los sistemas de gestión de residuos orgánicos. En particular, \textcite{galan-diazCarbonWaterFootprint2024} documentan cómo la implementación de sensores y automatización en granjas de \textit{BSF} permitió una reducción significativa de la huella de carbono y agua, gracias al control preciso de variables operativas como la temperatura, la humedad y el consumo energético. Este enfoque no solo mejora la eficiencia termodinámica del proceso, sino que habilita una trazabilidad digital rigurosa.

De forma complementaria, \textcite{jenkinsProcessingPoultryManure2025} evidencian que la integración de tecnologías inteligentes en el tratamiento de estiércol avícola contribuye a una disminución considerable en las emisiones de \ce{CO2} y \ce{N2O}, reforzando el papel de estas herramientas como catalizadores para la mitigación del cambio climático. En conjunto, estos estudios sugieren que la automatización no es solo una mejora productiva, sino un requisito habilitante para que los sistemas de tratamiento alcancen estándares de neutralidad en carbono y se integren efectivamente en estrategias de economía circular.

\section{Brechas tecnológicas y oportunidades de investigación}


A pesar de estos avances, la implementación operativa de sistemas inteligentes en la bioconversión con \textit{BSF} enfrenta barreras estructurales que limitan su escalabilidad y confiabilidad. La literatura identifica cinco brechas críticas que esta investigación busca abordar:
\begin{itemize}
\item Materialidad y Resiliencia del Hardware: Escasa disponibilidad de sensores de bajo costo capaces de mantener la precisión (calibración estable) en las condiciones agresivas del proceso (HR $>80\%$), lo que suele derivar en datos ruidosos o deriva instrumental \parencite{ahmedAgricultureClimateChange2020}.

\item Interoperabilidad: Falta de estandarización en protocolos de comunicación que impide la integración fluida entre dispositivos IoT heterogéneos y plataformas de gestión ambiental \parencite{vanIntegrationInternetofThingsSustainable2022}.

\item Limitaciones de los Modelos de Caja Negra: Los algoritmos de aprendizaje automático (AAA) puros, aunque potentes, carecen de restricciones físicas o biológicas explícitas. Esto los hace propensos al \textit{overfitting} y reduce su capacidad de generalización ante condiciones no vistas (como cambios de sustrato), a menos que se hibriden con modelos mecanísticos (EDO) \parencite{weichertReviewMachineLearning2019}.

\item Ausencia de Estándares MRV Dinámicos: Inexistencia de metodologías validadas para el Medición, Reporte y Verificación (MRV) en tiempo real, lo que excluye a la BSF de los mercados de carbono regulados \parencite{herrinCollateralDataQuality2021}.

\item Brecha de Implementación Rural: Dificultades para escalar soluciones dependientes de la nube en contextos con baja infraestructura de conectividad, exigiendo arquitecturas de procesamiento en el borde (\textit{Edge Computing}).
\end{itemize}

Estas brechas delinean una oportunidad estratégica para el desarrollo de una arquitectura híbrida compuesta por: (i) una red de sensores robustecida para monitoreo ambiental; (ii) modelos EDO para capturar la dinámica biológica base; (iii) algoritmos de corrección por AAA para adaptarse a la variabilidad real; y (iv) una interfaz de soporte a la decisión. Esta integración apunta a transformar la bioconversión de una caja negra biológica a un proceso transparente, tecnológicamente viable y verificable climáticamente \parencite{kamilaris2017review}.

\section{Brechas estructurales en el modelamiento matemático de la bioconversión}

Aunque el modelamiento matemático ha logrado avances significativos en la descripción bioenergética del crecimiento larval \parencite{padmanabhaComprehensiveDynamicGrowth2020,eriksenDynamicModellingFeed2022}, persisten brechas estructurales que limitan su validez ecológica y operativa. La principal limitación radica en el enfoque determinista y de \textit{estado estacionario} que predomina en la literatura: la mayoría de los modelos asumen condiciones ambientales homogéneas e idealizadas, ignorando los gradientes espaciales de temperatura y oxigenación que ocurren naturalmente en el lecho de sustrato. Esta simplificación impide predecir fenómenos emergentes críticos, como la formación de núcleos anaeróbicos generadores de \ce{CH4} y \ce{N2O}, cuya aparición es estocástica y dependiente de la microestructura del medio \parencite{bekkerImpactSubstrateMoisture2021a,rossiEstimatingDynamicsGreenhouse2024}.

Una segunda brecha crítica es la desconexión entre la teoría (modelo) y la observación (sensor). Los modelos vigentes operan en lazo abierto, incapaces de asimilar datos en tiempo real para corregir sus trayectorias ante perturbaciones no modeladas (ej. cambios súbitos en la composición del residuo). Esta rigidez contrasta con la naturaleza adaptativa de los sistemas biológicos, resultando en errores de predicción que aumentan con el tiempo de simulación. A esto se suma la falta de estandarización en las métricas de validación —con estudios que reportan emisiones en unidades funcionales incomparables—, lo que fragmenta el conocimiento y dificulta la consolidación de factores de emisión robustos para la industria \parencite{fuhrmannComprehensiveIndustryrelevantBlack2025}.

Estas limitaciones delinean una frontera de investigación clara: la necesidad de transitar de modelos puramente mecanísticos a arquitecturas de modelado híbrido. La integración de ecuaciones diferenciales (para la tendencia biológica) con algoritmos de aprendizaje automático alimentados por IoT (para la corrección de errores residuales) surge como la estrategia óptima para dotar al sistema de la capacidad adaptativa necesaria para operar en entornos agroindustriales reales \parencite{kobelskiModelbasedProcessOptimization2024a,liuOptimizingDataPipelines2023}.

\chapter{Objetivos de la Investigación}

\section{Objetivo General}

Desarrollar un sistema híbrido a escala de laboratorio —integrado por modelado mecanístico (EDO), sensores IoT de bajo costo y algoritmos de aprendizaje automático— para la estimación de la tasa de producción de \ce{CO2} y la detección de eventos de emisión de \ce{CH4} en la bioconversión de residuos con \textit{Hermetia illucens}, con el fin de generar métricas de desempeño ambiental aplicables a la gestión de residuos agroindustriales.

\section{Objetivos Específicos}

\begin{enumerate} \item \textbf{Instrumentación:} Configurar y calibrar un módulo de sensores NDIR de bajo costo para la medición de \ce{CO2} y \ce{CH4}, determinando su precisión y desviación (MAPE) frente a un instrumento de referencia bajo las condiciones de temperatura y humedad propias del cultivo de larvas.

\item \textbf{Modelado Matemático:} Adaptar y parametrizar un modelo de EDO basado en la bioenergética de \textit{Hermetia illucens} para estimar la tasa de respiración (\ce{CO2}), evaluando su capacidad predictiva en ventanas temporales durante un ciclo de bioconversión de 15 días.

\item \textbf{Integración Híbrida:} Desarrollar un algoritmo de estimación de estados mediante redes neuronales que infiera la biomasa larval a partir de datos operativos, integrando esta variable como entrada dinámica al modelo mecanístico para mejorar la precisión de la estimación de \ce{CO2}.

\item \textbf{Métricas Ambientales:} Validar la utilidad del sistema para la gestión ambiental mediante la detección de eventos de anaerobiosis (\ce{CH4}) y la cuantificación del carbono emitido, contrastando estos indicadores dinámicos frente a estimaciones estáticas tradicionales.
\end{enumerate}

\chapter{Marco Metodológico}

\section{Arquitectura del sistema híbrido}

El sistema implementa una arquitectura de modelado en cascada que integra técnicas de aprendizaje profundo con principios de bioenergética teórica. A diferencia de los enfoques de caja negra, esta propuesta utiliza una Red Neuronal Artificial (RNA) para estimar las variables de estado biológico (biomasa), las cuales alimentan posteriormente un modelo mecanístico basado en la teoría de Presupuesto Energético Dinámico (DEB), siguiendo los lineamientos de \textcite{eriksenDynamicModellingFeed2022}. 

La arquitectura se estructura en tres etapas secuenciales:

\textbf{1. Inferencia Biométrica (Red Neuronal):} Dado que la medición directa y continua de la biomasa larval es invasiva e impracticable en operación, se emplea una Red Neuronal, como estimador virtual. Este algoritmo recibe como entrada las series temporales de las variables ambientales (temperatura, humedad) y el histórico de emisiones, y tiene como objetivo reconstruir la curva de crecimiento de la biomasa húmeda ($W_w$) y determinar el estadio de desarrollo larval en tiempo real.

\textbf{2. Modelado Metabólico (Enfoque DEB):} Con la biomasa estimada, se ejecuta el modelo mecanístico DEB parametrizado para BSF. Este modelo desagrega matemáticamente el flujo de energía del organismo, calculando la tasa de respiración aeróbica obligada (producción de \ce{CO2}) en función de dos componentes: los costos de mantenimiento somático (proporcionales a la biomasa actual) y los costos de síntesis de nuevo tejido (proporcionales a la tasa de crecimiento). Esto genera una predicción teórica de cuánto \ce{CO2} \textit{debería} emitir el sistema si el metabolismo fuera puramente aeróbico y saludable.

\textbf{3. Diagnóstico por Contrastación (Residuo Operativo):} El nivel final compara la predicción teórica del modelo DEB (el "deber ser"\ biológico) con la medición real de los sensores NDIR. La discrepancia entre ambos valores actúa como un indicador de desempeño:\begin{itemize}\item Una coincidencia alta valida la eficiencia del proceso.\item Una divergencia significativa (ej. sensor $<$ modelo) sugiere inhibición metabólica o problemas en el sistema.\item La presencia simultánea de \ce{CH4} (medido) confirma desviaciones hacia rutas anaerobias no contempladas en el modelo DEB, gatillando alertas de gestión.\end{itemize}

\section{Ubicación y Contexto Experimental}

El componente experimental se desarrollará en la sede Palmira de la Universidad Nacional de Colombia, ver Figura \ref{fig:localizacion}). La zona se caracteriza por corresponder a un ecosistema de bosque tropical seco, con una temperatura media anual de $24^\circ C$, condiciones que favorecen el desarrollo natural de BSF. El montaje se realizará en condiciones de laboratorio controlado (mesocosmos) para garantizar la reproducibilidad de los datos de calibración.

\begin{figure}[h]
    \centering
    \includegraphics[width=0.8\textwidth]{00Figuras/localizacion.jpg} % Asegúrate de que la ruta sea correcta
    \caption{Ubicación del proyecto en Palmira, Valle del Cauca, Colombia.}
    \label{fig:localizacion}
\end{figure}

\section{Enfoque y Tipo de Investigación}

Este estudio se enmarca dentro de una investigación aplicada, experimental, tecnocientífica e interdisciplinaria, orientada al desarrollo y validación de un sistema híbrido para la caracterización dinámica de emisiones de gases de efecto invernadero (GEI).

El enfoque es cuantitativo, fundamentado en la recolección sistemática de series de tiempo mediante sensores IoT, el modelado determinista (EDO) y el análisis estadístico predictivo (Machine Learning). La metodología se estructura en cuatro fases secuenciales e integradas que van desde la instrumentación física hasta la validación de la utilidad ambiental del sistema.

\section{Diseño Experimental y Desarrollo (Materiales y Métodos)}

\subsection{Fase I: Diseño e Implementación del Sistema IoT y Aseguramiento de Datos}

El objetivo de esta fase es garantizar la calidad del dato de entrada. Dada la naturaleza agresiva del entorno de bioconversión (alta humedad relativa $>80\%$), se implementarán protocolos de robustez instrumental.

\textbf{Actividades Clave:}
\begin{itemize}
    \item \textbf{Implementación de Hardware:} Diseño de un nodo de sensores basado en microcontroladores (tipo ESP32), integrando sensores NDIR para \ce{CO2} y semiconductores/electroquímicos para \ce{CH4}.
    \item \textbf{Protocolos de Calibración:} Ejecución de curvas de calibración frente a concentraciones patrón y desarrollo de algoritmos de \textbf{compensación cruzada} para corregir la deriva del sensor causada por las fluctuaciones de temperatura y humedad inherentes al proceso exotérmico.
    \item \textbf{Captura de Datos:} Despliegue de la red IoT para la recolección de datos en ciclos de bioconversión (14-20 días), generando el \textit{dataset} experimental necesario para el entrenamiento de los modelos.
\end{itemize}

\subsection{Fase II: Formulación y Parametrización del Modelo Mecanístico (Línea Base)}

Se desarrollará un modelo matemático basado en Ecuaciones Diferenciales Ordinarias (EDO) siguiendo el marco del \textit{Dynamic Energy Budget} (DEB) simplificado para BSF. La función de este módulo es generar una predicción de referencia de la respiración del sistema bajo condiciones ideales.

\textbf{Enfoque de Modelado Bioenergético:}
\begin{itemize}
    \item \textbf{Dinámica del \ce{CO2} ($J_{CO2}$):} El modelo describe la producción de dióxido de carbono como la suma de dos procesos fisiológicos acoplados: el costo energético de mantenimiento (proporcional a la biomasa actual) y el costo de biosíntesis (proporcional a la tasa de crecimiento), según la lógica de \textcite{eriksenDynamicModellingFeed2022}.
    \item \textbf{Supuesto de Aerobiosis Estricta:} Se asume que el sistema opera con suministro de oxígeno no limitante. Por consiguiente, la generación de \ce{CH4} se excluye de las ecuaciones de estado deterministas, considerándose una variable de perturbación estocástica indicativa de fallos operativos (zonas anóxicas).
    \item \textbf{Salida del Modelo:} El resultado es una curva teórica de respiración ($Y_{DEB}$) que representa el \textit{comportamiento esperado}. Esta señal servirá de línea base para calcular los residuales frente a las mediciones reales en la etapa de validación.\end{itemize}

\subsection{Fase III: Desarrollo del Sistema Híbrido (Validación Cruzada)}

En esta fase se integra la información experimental determinista (biomasa medida) con la observación estocástica (sensores) para validar el desempeño del proceso. La arquitectura resuelve el conflicto de incertidumbre mediante la siguiente lógica de jerarquización:

\textbf{Estrategia de Validación:}

\begin{itemize}\item \textbf{Entrada de Biomasa :} Los datos experimentales de biomasa larval ($W_{real}$), obtenidos mediante pesaje periódico, constituyen la variable independiente de entrada. Dado que estos datos son discretos (ej. mediciones cada 48h), se ajustará una curva de crecimiento (función logística o Gompertz) para proveer un valor continuo de biomasa ($W(t)$) al modelo en cada instante de tiempo.\item \textbf{Generación de la Expectativa Ideal ($Y_{modelo}$):} El modelo EDO/DEB recibe el valor de $W(t)$ y calcula la tasa de respiración teórica esperada bajo condiciones de aerobiosis estricta. Este valor representa la línea base de eficiencia máxima: ¿Cuánto \ce{CO2} debería generar esta cantidad exacta de larvas si estuvieran totalmente sanas y el medio fuera 100\% aeróbico?.

\item \textbf{Análisis de Discrepancia (El aporte del Sensor):} Se contrasta la curva teórica con la señal del sensor ($Y_{sensor}$). Se asume que:
\begin{equation}
    Y_{sensor} = Y_{modelo} + \delta_{bio} + \epsilon_{inst}
\end{equation}
Donde $\epsilon_{inst}$ es el error instrumental del sensor (asumido acotado por calibración H1) y $\delta_{bio}$ es la \textit{desviación biológica}.

\item \textbf{Criterio de Decisión:} Si la diferencia $|Y_{sensor} - Y_{modelo}|$ supera el umbral de error del sensor y correlaciona con la presencia de \ce{CH4} (medido), el sistema atribuye la causa a una ineficiencia del proceso (anaerobiosis), validando así la utilidad del modelo como herramienta de referencia para la gestión ambiental.
\end{itemize}

\subsection{Fase IV: Evaluación del Desempeño del Prototipo}

Esta fase final tiene por objeto cuantificar la diferencia práctica entre medir el proceso continuamente frente a estimarlo teóricamente.

\textbf{Estrategia de Evaluación:} \begin{itemize} \item \textbf{Detección de Eventos (Validación de H4):} Se verificará si el sistema fue capaz de registrar picos de \ce{CH4} en momentos donde el modelo teórico predecía normalidad. El éxito se medirá de forma binaria: ¿Detectó el sensor el evento anómalo? (Sí/No).

\item \textbf{Cálculo de la Discrepancia de Masa:} Se comparará la masa total de carbono emitida según el sensor (integrando el área bajo la curva medida) contra la masa estimada por el modelo teórico.
\begin{equation}
    \Delta_{Total} = \int_{0}^{t} (Y_{sensor} - Y_{modelo}) \, dt
\end{equation}
Este valor $\Delta_{Total}$ representará la cantidad de gas que el método tradicional "no vio" (o sobreestimó), sirviendo como indicador simple de la utilidad del monitoreo en tiempo real.
\end{itemize}

\section{Arquitectura Tecnológica del Sistema}

El sistema sigue una arquitectura modular de tres capas, diseñada bajo el paradigma IoT para garantizar la integridad, disponibilidad y trazabilidad del dato científico (Ver Figura \ref{fig:arquitectura}):

\subsection{Capa de Percepción (Edge Computing)} Compuesta por el nodo sensor (Hardware basado en ESP32). Su función principal es la digitalización de variables físicas. Para garantizar la robustez del dato experimental, esta capa implementa: \begin{itemize} \item Adquisición y Filtrado: Lectura de sensores y aplicación de filtros digitales (media móvil) para atenuar el ruido eléctrico. \item Resiliencia: Almacenamiento temporal en memoria local (SD/Flash) para evitar pérdida de datos ante fallos de conectividad (buffer circular). \item Telemetría: Transmisión asíncrona de la trama de datos mediante protocolo ligero MQTT hacia el servidor de procesamiento. \end{itemize}

\subsection{Capa de Procesamiento (Core/Cloud)} Módulo de software centralizado donde reside la lógica científica del sistema. \begin{itemize} \item Gestor de Datos (Ingesta): Recibe la telemetría, realiza una validación de rangos (sanitización de datos) y la almacena en una base de datos relacional para garantizar la persistencia histórica. \item Motor Híbrido de Estimación: Ejecuta secuencialmente el modelo matemático: primero resuelve la cinética base (EDO) y posteriormente invoca el algoritmo de aprendizaje automático para corregir el error residual, generando la "Variable Estimada"\ final. \end{itemize}

\subsection{Capa de Aplicación (Interfaz de Decisión)} Interfaz de visualización orientada al usuario final y al soporte de decisiones operativas. Sus funcionalidades clave incluyen: \begin{itemize} \item Monitorización Dinámica: Visualización de las curvas de tendencia (Real vs. Estimada) en tiempo real. \item Auditoría Ambiental: Módulo de descarga de reportes consolidados del ciclo para procesos de MRV. \item Gestión de Alertas: Semáforo de riesgo que notifica visualmente la detección de patrones compatibles con anaerobiosis incipiente. \end{itemize}

\begin{figure}[h!] \centering \includegraphics[width=1.0\textwidth]{00Figuras/DG_FLUJO_GENERAL.png} \caption{Diagrama de la arquitectura funcional del sistema de monitoreo y predicción híbrido, detallando el flujo de datos desde la percepción física hasta la toma de decisiones.} \label{fig:arquitectura} \end{figure}

\section{Criterios de Validación y Métricas de Éxito}Para determinar el cumplimiento de los objetivos e hipótesis, se establecen los siguientes umbrales de desempeño, alineados con el alcance de un prototipo experimental:\begin{itemize}\item \textbf{Robustez Instrumental (H1):} Se aceptará una correlación de Pearson positiva ($R > 0.85$) entre la señal de los sensores NDIR y el instrumento de referencia, con un Error Porcentual Absoluto Medio ($MAPE$) $\leq 15\%$, un margen tolerado considerando la deriva instrumental en ambientes de alta humedad relativa ($>80\%$).\item \textbf{Precisión del Modelo Híbrido (H2 y H3):} El sistema integrado (Estimación de Biomasa IA $\to$ Modelo DEB) se considerará exitoso si logra predecir la tasa de evolución de \ce{CO2} con un Coeficiente de Determinación ($R^2$) $\geq 0.70$ respecto a los datos observados. Además, se busca que el error de predicción (RMSE) sea inferior al obtenido mediante un modelo de crecimiento lineal estático, demostrando el valor de capturar la dinámica biológica.

\item \textbf{Sensibilidad a Anomalías (H4):} Se validará la capacidad del sistema para actuar como alerta temprana de anaerobiosis. La clasificación de eventos (Presencia/Ausencia de picos de \ce{CH4}) deberá alcanzar una \textbf{Sensibilidad (\textit{Recall})} $> 80\%$. Se prioriza esta métrica sobre la precisión pura, para el apoyo en la gestión ambiental.

\item \textbf{Valor de Gestión (Objetivo General):} El sistema deberá entregar un \textbf{Inventario Dinámico de Carbono}, cuantificando la discrepancia de masa ($\Delta_{Total}$) entre la emisión real acumulada y la estimación teórica. El criterio de éxito será la capacidad del sistema para explicitar esta diferencia (en gramos de C) al final del ciclo, aportando un dato de incertidumbre reducida para el reporte ambiental.
\end{itemize}
\section{Alineación de Objetivos y Estrategia Experimental}

A modo de síntesis, la Tabla~\ref{tab:resumen_metodologia} presenta la matriz de coherencia que articula los objetivos específicos con sus respectivas fases de ejecución y criterios de validación. Esta estructura evidencia que el diseño metodológico no opera como un conjunto de técnicas aisladas, sino como una arquitectura en cascada integrada. En este esquema, la \textbf{precisión instrumental} (Fase I) alimenta la \textbf{inferencia de estados biológicos} mediante inteligencia artificial (Fase III), lo cual a su vez habilita la \textbf{fidelidad bioenergética} del modelo mecanístico (Fase II). Este encadenamiento lógico asegura que cada etapa contribuya de forma directa y cuantificable a la generación de métricas de desempeño ambiental (Fase IV), garantizando la trazabilidad científica desde la captura de la señal física hasta la toma de decisiones operativas.

\begin{table}[h]
\centering
\caption{Matriz de Coherencia Metodológica y Criterios de Éxito}
\label{tab:resumen_metodologia}
\renewcommand{\arraystretch}{1.3} % Para dar un poco más de aire a las filas
\begin{tabular}{|p{3cm}|p{3.5cm}|p{4.5cm}|p{3.5cm}|}
\hline
\textbf{Objetivo Específico} & \textbf{Método / Fase} & \textbf{Entregable Clave} & \textbf{Criterio de Éxito} \\ \hline
1. Instrumentación & Fase I: Calibración y Despliegue & Módulo de sensores validado en ambiente saturado & $R > 0.85$ y MAPE $< 15\%$ (vs. referencia) \\ \hline
2. Modelo EDO & Fase II: Parametrización Bioenergética & Ecuaciones de tasa respiratoria parametrizadas para \textit{H. illucens} & Capacidad de replicar la tendencia base ($R^2 > 0.70$) \\ \hline
3. Integración Híbrida & Fase III: Estimación de Estados (IA $\to$ DEB) & Algoritmo de estimación de biomasa y predicción dinámica de \ce{CO2} & Error (RMSE) inferior al de un modelo de crecimiento lineal estático \\ \hline
4. Métricas Ambientales & Fase IV: Validación Operativa & Inventario Dinámico de Carbono y sistema de alertas & Recall $> 80\%$ en detección de eventos de Anaerobiosis (\ce{CH4}) \\ \hline
\end{tabular}
\end{table}
\section{Análisis Crítico: Limitaciones, Riesgos y Desafíos de Implementación}

Si bien la integración de tecnologías digitales promete optimizar la bioconversión, este estudio adopta una postura \textit{tecno-realista}, reconociendo que la digitalización no es neutral ni está exenta de externalidades. En respuesta a la necesidad de un análisis riguroso sobre la viabilidad técnica y socioambiental, a continuación se detallan las tensiones éticas, ambientales, epistémicas y operativas inherentes a la propuesta, junto con las estrategias de mitigación integradas en su diseño.

\subsection{Tensión socioecológica: impacto de la materialidad digital}
Existe una tensión entre la sostenibilidad del proceso biológico (bioconversión) y la huella ecológica asociada a la infraestructura digital. La implementación de redes IoT conlleva impactos en la extracción de minerales críticos, el consumo energético y la generación futura de residuos electrónicos (\textit{e-waste}).

\begin{itemize}
    \item \textbf{Riesgo (Degradación prematura):} La exposición prolongada a condiciones agresivas —humedad relativa $>80\%$, amoniaco y compuestos volátiles corrosivos— puede acortar drásticamente la vida útil de los sensores, incrementando la frecuencia de reemplazos y la generación de residuos tecnológicos.
    \item \textbf{Mitigación (Diseño para mantenimiento):} El sistema se basa en una arquitectura modular que permite reemplazar únicamente el elemento sensible (sensor), preservando la unidad de procesamiento (ESP32) y comunicación. Esta estrategia, combinada con el uso de microcontroladores de bajo consumo, reduce tanto la huella operativa como la tasa de generación de \textit{e-waste} \parencite{dharCarbonImpactArtificial2020}.
\end{itemize}

\subsection{Riesgos epistémicos: opacidad algorítmica y sobreajuste}
La dependencia de algoritmos de aprendizaje automático introduce incertidumbres sobre la interpretabilidad y la capacidad de generalización del modelo en entornos no estacionarios.



%[Image of data processing flowchart]


\begin{itemize}
    \item \textbf{Limitación (Sobreajuste a condiciones de laboratorio):} Los modelos de aprendizaje automático puros pueden ajustarse excesivamente a las condiciones controladas del prototipo, perdiendo robustez al enfrentarse a la variabilidad de sustratos y microclimas en operación real.
    \item \textbf{Estrategia (Modelado híbrido con restricciones físicas):} Se adopta un enfoque de inteligencia artificial explicable (XAI) mediante la integración con un modelo mecanístico (EDO) basado en el marco del \textit{Dynamic Energy Budget} (DEB). Esta integración no solo corrige el error residual, sino que impone restricciones biológicas en la función de pérdida del algoritmo, evitando predicciones no realistas (e.g., tasas de crecimiento negativas o emisiones superiores al carbono disponible) y manteniendo la interpretabilidad para el operario.
\end{itemize}

\subsection{Retos de escalamiento y brecha digital en contextos rurales}
La transición del prototipo (TRL 4) a una implementación a escala (TRL 7–9) enfrenta barreras estructurales en contextos latinoamericanos, donde la conectividad y el acceso a infraestructura digital son desiguales.



%[Image of edge computing architecture diagram]


\begin{itemize}
    \item \textbf{Desafío de conectividad:} Las arquitecturas centralizadas en la nube excluyen a productores en zonas con baja o nula cobertura móvil ("zonas grises"), potencialmente exacerbando inequidades tecnológicas.
    \item \textbf{Oportunidad (Computación en el borde):} La arquitectura propuesta descentraliza la inteligencia mediante \textit{edge computing}, ejecutando los algoritmos críticos directamente en el nodo IoT. Esto permite operación autónoma y sincronización asíncrona de datos, garantizando funcionalidad incluso sin conexión continua.
    \item \textbf{Accesibilidad económica:} El escalamiento depende de que el costo del sistema sea marginal respecto al valor del producto. El uso de plataformas embebidas de bajo costo (ESP32) y sensores comerciales busca democratizar el acceso para pequeños y medianos productores agroindustriales.
\end{itemize}

\subsection{Bioseguridad y gobernanza de datos}
El proyecto se alinea con principios de innovación responsable y transparencia climática, cumpliendo con marcos regulatorios emergentes.

\begin{itemize}
    \item \textbf{Bioseguridad:} Aunque BSF no es vector de patógenos humanos, el diseño del sistema refuerza los protocolos de contención biológica. Mediante la detección temprana de anomalías en los parámetros ambientales (fugas o cambios térmicos súbitos), el sistema genera alertas críticas para la intervención manual inmediata, apoyando el cumplimiento de normativas fitosanitarias (ICA/UE) sin depender exclusivamente de barreras pasivas.
    
    \item \textbf{Integridad MRV:} Los registros de emisiones se gestionan bajo principios de trazabilidad digital y sellado de tiempo (\textit{timestamping}), garantizando que los datos sean auditables y consistentes con los requisitos de Medición, Reporte y Verificación (MRV) de estándares internacionales (IPCC, Verra). Esto mitiga el riesgo de \textit{greenwashing digital} y sienta las bases técnicas para la futura certificación de créditos de carbono.
\end{itemize}

\section{Fase 4: Integración para la gestión socio-ambiental y operativa}

Para trascender el alcance puramente técnico y responder a la necesidad de una gestión integral, el sistema no se limita a la generación de datos, sino que integra un esquema de \textit{respuesta operativa} basado en la información predicha. Esta fase conecta la salida del algoritmo con la toma de decisiones humanas bajo el siguiente protocolo:

\subsection{Protocolo de respuesta ante eventos de anaerobiosis}
La predicción de picos de \ce{CH4} (indicador de anaerobiosis) activará alertas tempranas en el dashboard, diseñadas no solo como métricas técnicas, sino como instrucciones operativas claras para el productor (ej. Riesgo de metano alto: Se requiere volteo manual o aumento de aireación).
\begin{itemize}
    \item \textbf{Nivel Verde (Normal):} Emisiones dentro del rango aeróbico estándar. El sistema valida la acumulación de créditos de carbono.
    \item \textbf{Nivel Ámbar (Preventivo):} Tendencia al alza en la relación \ce{CO2}/\ce{O2}. Sugiere revisión de humedad.
    \item \textbf{Nivel Rojo (Crítico):} Detección de picos de \ce{CH4}. Se instruye intervención inmediata (volteo) para reoxigenar el sustrato, mitigando el impacto ambiental en tiempo real.
\end{itemize}

Este mecanismo garantiza que la predicción no sea un dato pasivo, sino un activador de control de procesos que reduce directamente la huella de carbono del tratamiento.

\chapter{Estimación de Costos y Presupuesto}

\section{Resumen Financiero del Proyecto}
El desarrollo de esta investigación implica la implementación de una arquitectura híbrida de hardware y software para la validación de métricas ambientales. La estructura de costos se ha diseñado bajo el principio de eficiencia, priorizando el desarrollo de capital intelectual y la validación experimental sobre la adquisición de activos fijos costosos.

El presupuesto global estimado asciende a $267$ millones de pesos colombianos (COP), ejecutados a lo largo de 4 años. Este monto cubre desde la fase de diseño conceptual hasta la validación en campo y la divulgación científica.
\section{Desglose de Costos de Desarrollo Tecnológico}

A continuación se detalla la inversión específica requerida para la construcción del prototipo funcional (TRL 5-6), diferenciando componentes físicos y servicios digitales.

\begin{table}[H]
    \centering
    \begin{tabular}{|p{7cm}|c|c|}
        \hline
        \textbf{Componente Tecnológico} & \textbf{Unidades} & \textbf{Costo (MM COP)} \\ \hline
        \multicolumn{3}{|l|}{\textit{\textbf{1. Hardware y Sensórica (Nodo IoT)}}} \\ \hline
        Sensores de \ce{CO2} (NDIR Grado Industrial) & 3 & 0.6 \\ \hline
        Sensores de Gases Traza (\ce{CH4}, VOCs) & 3 & 0.72 \\ \hline
        Sensores Ambientales (Temp/Humedad de alta precisión) & 3 & 1.2 \\ \hline
        Sensores Auxiliares (Luz, pH, NPK) & 3 & 0.6 \\ \hline
        Unidades de Procesamiento (ESP32 + LoRa) & 5 & 4.0 \\ \hline
        Manufactura de PCB y Carcasas IP65 & 3 & 0.8 \\ \hline
        \multicolumn{3}{|l|}{\textit{\textbf{2. Software e Infraestructura de Datos}}} \\ \hline
        Desarrollo de Algoritmos Híbridos (Licenciamiento) & Global & 4.0 \\ \hline
        Interfaz de Visualización (Dashboard Web) & Global & 2.0 \\ \hline
        Infraestructura Cloud (Hosting + DB - 12 meses) & Global & 2.4 \\ \hline
        \multicolumn{3}{|l|}{\textit{\textbf{3. Validación y Transferencia}}} \\ \hline
        Equipos Patrón de Referencia (Alquiler/Uso) & Global & 2.0 \\ \hline
        Insumos para Ensayos Experimentales & Global & 2.8 \\ \hline
        Kit de Comunicaciones (Gateway LoRaWAN + Antenas) & 1 & 1.8 \\ \hline
        \multicolumn{3}{|l|}{\textit{\textbf{4. Servicios Especializados}}} \\ \hline
        Consultoría Técnica Externa & Global & 6.0 \\ \hline
        Capacitación y Formación Especializada & Global & 3.2 \\ \hline
        \textbf{SUBTOTAL DESARROLLO TECNOLÓGICO} & & \textbf{\$ 32.12} \\ \hline
    \end{tabular}
    \caption{Estimación de costos directos de desarrollo e implementación.}
    \label{tab:costos_desarrollo}
\end{table}

\section{Cronograma de Ejecución Financiera}

La siguiente tabla detalla la distribución de recursos por actividad estratégica, indicando el estado actual de avance y la valoración económica de los entregables (recursos humanos, técnicos y operativos).

\begin{longtable}{|p{5cm}|p{6cm}|c|c|}
    \hline
    \textbf{Actividad} & \textbf{Descripción del Rubro} & \textbf{\% Avance} & \textbf{Costo (MM COP)*} \\
    \hline
    \endfirsthead
    \multicolumn{4}{c}{\textit{Continúa de la página anterior}} \\
    \hline
    \textbf{Actividad} & \textbf{Descripción} & \textbf{\% Avance} & \textbf{Costo (MM COP)} \\
    \hline
    \endhead
    \hline
    \multicolumn{4}{r}{\textit{*MM COP: Millones de Pesos Colombianos}} \\
    \endfoot
    \hline
    \endlastfoot
    
    Revisión bibliográfica & Acceso a bases de datos y análisis del estado del arte. & 50\% & 6.0 \\
    \hline
    Infraestructura Computacional & Uso de equipos de cómputo, licencias de software y simulación. & 70\% & 8.0 \\
    \hline
    Ingeniería de Software & Desarrollo de algoritmos de backend y procesamiento (4 años). & 45\% & 20.0 \\
    \hline
    Ingeniería de Hardware & Diseño PCB, integración de sensores y manufactura de nodos. & 100\% & 20.0 \\
    \hline
    \textbf{Estudio Doctoral} & \textbf{Dedicación del investigador (Matrícula + Manutención 4 años).} & 45\% & 80.0 \\
    \hline
    Validación Experimental & Logística de campo, insumos biológicos (larvas) y sustratos. & 0\% & 16.0 \\
    \hline
    Desarrollo de IA & Entrenamiento de modelos y horas de procesamiento en GPU. & 10\% & 24.0 \\
    \hline
    Sistema MRV & Diseño de la plataforma de reporte y validación ambiental. & 60\% & 12.0 \\
    \hline
    Servicios en la Nube & Costos operativos de almacenamiento y despliegue (AWS/Azure). & 50\% & 10.0 \\
    \hline
    Divulgación Científica & Publicación en revistas y asistencia a congresos. & 25\% & 8.0 \\
    \hline
    Gestión del Proyecto & Administración, papelería y costos indirectos. & 60\% & 8.0 \\
    \hline
    Metrología y Calibración & Servicios de laboratorio certificado para validación de sensores. & 40\% & 12.0 \\
    \hline
    Escritura y Edición & Elaboración de documentos, traducción y corrección de estilo. & 20\% & 8.0 \\
    \hline 
    \textbf{TOTAL ESTIMADO} &  & \textbf{-} & \textbf{\$ 232.0} \\
    \hline
\end{longtable}

\section{Estrategia de Financiamiento}

La viabilidad financiera del proyecto se sustenta en un modelo mixto de financiación:

\begin{itemize}
    \item \textbf{Aporte Institucional (UNAL):} Cubre principalmente los costos de capital humano (dirección de tesis), acceso a laboratorios, infraestructura física y recursos bibliográficos, valorados como contrapartida en especie.
    \item \textbf{Recursos Propios:} El investigador asume los costos directos de adquisición de componentes electrónicos importados, insumos de experimentación y gastos operativos de campo no cubiertos por convocatorias institucionales.
\end{itemize}
\section{Cronograma de Actividades}

El plan de trabajo se estructura en un horizonte temporal de 42 meses (agosto 2023 – enero 2027), organizado mediante una ejecución paralela de actividades transversales y fases de desarrollo tecnológico incremental. La distribución temporal, detallada en la Figura \ref{fig:cronograma}, obedece a la siguiente lógica operativa:

\begin{description} \item[Actividades Transversales (2023-2027):] Se contempla la \textbf{Revisión Bibliográfica} y el desarrollo de la \textbf{Tesis Doctoral} como procesos continuos durante todo el ciclo del proyecto. Esto garantiza que tanto el marco teórico como la escritura del documento final se actualicen constantemente con los hallazgos experimentales.

\item[Fase 1: Infraestructura y Fundamentación (2023-2024):] 
Inicia con el despliegue temprano de la \textbf{Infraestructura en la Nube} (finales de 2023), asegurando el entorno digital antes de la llegada de los datos físicos. Simultáneamente, se consolidan las bases teóricas del estudio.

\item[Fase 2: Ingeniería Hard/Soft (2024-2026):] 
Es el núcleo de desarrollo técnico. Inicia en el segundo semestre de 2024 con la \textbf{Ingeniería de Software} y \textbf{Hardware}. Nótese que el diseño de sensores IoT se extiende hasta mediados de 2025, mientras que la ingeniería de hardware continúa hasta inicios de 2026 para permitir ajustes iterativos post-validación.

\item[Fase 3: Inteligencia Artificial y Validación (2025-2026):] 
A partir de 2025, con los primeros prototipos funcionales, inicia el \textbf{Desarrollo de IA} (que corre paralelo durante dos años) y la \textbf{Validación Experimental} (mediados de 2025 a mediados de 2026). Esta superposición es intencional para permitir el entrenamiento de algoritmos con datos reales a medida que se generan.

\item[Fase 4: Consolidación y Divulgación (2026-2027):] 
La etapa final se centra en la implementación del sistema \textbf{MRV} (segundo semestre de 2026), la producción de \textbf{Publicaciones Científicas} y el cierre administrativo y académico del doctorado en enero de 2027.
\end{description}

\begin{sidewaysfigure} 
    \centering 
    \includegraphics[width=1.0\textwidth]{00Figuras/gant.png} 
    \caption{Cronograma general de investigación y desarrollo tecnológico (2023-2027).} \label{fig:cronograma} 
\end{sidewaysfigure}

\section{Retorno de la Inversión (Justificación)}
Más allá del aporte académico, la inversión en este sistema genera valor tangible al reducir la incertidumbre en la medición de emisiones. Un sistema MRV validado habilita el acceso a mercados de carbono, cuyo potencial económico para una planta de bioconversión de 10 toneladas/día podría superar los costos de implementación tecnológica en menos de 2 años de operación, demostrando la sostenibilidad financiera de la propuesta a escala industrial.


\chapter{Resultados Esperados}

Conforme al alcance definido y la ventana de ejecución experimental, el proyecto se orienta a la validación de un sistema tecnológico funcional a escala de laboratorio. El propósito central no es solo la medición instrumental, sino la demostración de que una arquitectura híbrida (Sensores + IA + Bioenergética) puede reducir la incertidumbre en el monitoreo de procesos biológicos complejos, estableciendo bases técnicas para la gestión basada en datos y la reducción de la huella de carbono en la agroindustria.A continuación, se detallan los productos técnicos, científicos y ambientales comprometidos:

\section{Prototipo Funcional de Monitoreo IoT}Se entregará un sistema de adquisición de datos operativo y caracterizado, diseñado para mantener la integridad de la señal en el entorno agresivo de la bioconversión (alta humedad relativa y presencia de gases corrosivos).\begin{itemize}\item \textbf{Entregable:} Nodo de sensores integrado (basado en ESP32 u homólogo) capaz de registrar y transmitir telemetría de \ce{CO2}, \ce{CH4}, temperatura y humedad con alta resolución temporal. El hardware incluirá estrategias de preacondicionamiento de señal para mitigar el ruido electrónico inherente al entorno.\item \textbf{Criterio de Éxito:} Continuidad operativa documentada durante ciclos completos de 14 días, con una tasa de pérdida de paquetes de datos $< 5\%$. Se validará la calidad del dato mediante una correlación de Pearson ($R > 0.85$) y un error medio ($MAPE \leq 15\%$) frente a instrumentos de referencia, demostrando la viabilidad del bajo costo pese a la deriva instrumental.\end{itemize}

\section{Modelo Híbrido de Estimación de Estados}Se desarrollará un motor de cómputo que integre la inteligencia artificial con la teoría biológica para superar la imposibilidad de medir la biomasa larval en tiempo real.\begin{itemize}\item \textbf{Entregable:} Un algoritmo validado de Arquitectura en Cascada, compuesto por:\begin{enumerate}\item Un estimador basado en Redes Neuronales que infiera la biomasa larval ($W_w$) a partir de las variables ambientales históricas.\item Un modelo mecanístico (EDO-DEB) que utilice dicha biomasa para predecir la tasa de respiración (\ce{CO2}) esperada.\end{enumerate}\item \textbf{Criterio de Éxito:} El sistema híbrido deberá superar la precisión de una línea base estática (modelo de crecimiento lineal), logrando un coeficiente de determinación ($R^2 \geq 0.70$) en la predicción de la tendencia de \ce{CO2}, demostrando que el sistema captura efectivamente la dinámica de crecimiento de BSF.\end{itemize}

\section{Herramienta de Diagnóstico Ambiental}Se entregará una metodología validada para transformar los datos crudos en indicadores de gestión, útiles para la toma de decisiones operativas y el reporte de sostenibilidad.\begin{itemize}\item \textbf{Entregable (Sistema de Alertas):} Un módulo de detección lógica capaz de identificar eventos de Anaerobiosis mediante la lectura directa de \ce{CH4} y el análisis de divergencia entre el modelo y el sensor.\item \textbf{Entregable (Inventario Dinámico):} Un reporte automatizado que cuantifique la masa total de carbono emitida ($C_{total}$), explicitando la discrepancia ($\Delta$) entre la medición real y la estimación teórica tradicional.\item \textbf{Criterio de Éxito:} Una sensibilidad (\textit{Recall}) $> 80\%$ en la clasificación correcta de eventos de pico de metano, asegurando que el sistema actúe como una barrera efectiva contra emisiones fugitivas no reportadas.\end{itemize}

\section{Activo de Información Científica}Como resultado transversal, se consolidará un Dataset Experimental Estructurado sobre la dinámica de emisiones en la cría de BSF.\begin{itemize}\item \textbf{Entregable:} Base de datos curada y etiquetada que correlaciona variables ambientales, crecimiento de biomasa (ground truth) y emisiones de GEI (\ce{CO2}/\ce{CH4}) en alta frecuencia. Este activo servirá como insumo para futuras investigaciones o para el re-entrenamiento de algoritmos más complejos.\end{itemize}


\section{Interfaz de Visualización y Soporte Operativo}Para garantizar que la información técnica sea accesible para la toma de decisiones, se implementará una capa de visualización de datos de bajo costo computacional.\begin{itemize}\item \textbf{Entregable:} Un Dashboard de Control (basado en plataformas \textit{open-source} como Grafana, Streamlit o ThingsBoard) que despliegue en tiempo real:\begin{enumerate}\item Las curvas de concentración de \ce{CO2} y \ce{CH4}.\item El estado semafórico del sistema (Verde: Eficiente / Rojo: Alerta de Anaerobiosis).\item El valor acumulado de emisiones estimada vs. real.\end{enumerate}\item \textbf{Criterio de Éxito:} La interfaz deberá actualizarse automáticamente con una latencia $< 1$ minuto respecto a la transmisión del sensor, permitiendo la interpretación visual inmediata de las anomalías por parte de un operador no experto.\end{itemize}

\section{Producción Intelectual y Divulgación}Como validación externa de la calidad científica del desarrollo.\begin{itemize}\item \textbf{Entregable:} Un manuscrito científico (tipo \textit{Conference Paper} o artículo de revista indexada) que documente la metodología de calibración de sensores \textit{low-cost} en ambientes de alta humedad o la validación del modelo híbrido para BSF.\item \textbf{Criterio de Éxito:} Sometimiento o aceptación del trabajo en un evento académico o revista del área de Ingeniería Ambiental, Electrónica o Bioeconomía.\end{itemize}

%\include{01Seccion01}
%\include{02Seccion02}
%\include{03Seccion03}
%\include{04Seccion04}
%\include{05Seccion05}
%\include{06Seccion06}
%\include{07Seccion07}

% Inicio del apéndice o anexos
%\begin{appendix}
%\include{08Apendice01}%
%\end{appendix}

% Permite visualizar la bibliografía en la tabla de contenido
\addcontentsline{toc}{chapter}{Referencias Bibliográficas} % Literatura Citada

% Formato de la bibliografía
\pagestyle{plain} % Estilo de página para la bibliografía
\renewcommand{\bibfont}{\scriptsize} % Cambia el tamaño de la fuente de la bibliografía

% Nombre del documento donde se almacenan las referencias
\printbibliography
%\nocite{*} % Incluye todas las referencias, incluso las no citadas

\end{document}